%=======================================================================
%  CHAPTER 4 — SOLVED PROBLEMS
%  Every worked answer follows the Insights method, checked independently.
%  Where the book is wrong the notes do it correctly and say so.
%=======================================================================
\section{Knowledge Representation, Inference and Reasoning \textcolor{sub}{\large -- Solved Problems}}

{\color{sub}\footnotesize \mk{The biggest companion in the subject.}
33 papers set a resolution refutation, 15 set a Bayes calculation, 4 set a CNF conversion
and 1 sets a truth-table proof. Every one is worked below under the question
\textbf{exactly as its paper prints it}. Renamings of a puzzle already worked get the
paper's own statement plus what changes, so nothing is left to guess and no proof is
copied out twice.}

\lead{The five procedures everything here is built on}
\begin{itemize}
  \item \textbf{P1 Truth table} $-$ $2^n$ rows for $n$ symbols; a \textbf{tautology} is a
        column of all T.
  \item \textbf{P2 CNF} $-$ eliminate $\leftrightarrow$, eliminate $\rightarrow$, move
        $\neg$ in, standardise apart, skolemize, drop $\forall$, distribute $\vee$ over
        $\wedge$, split.
  \item \textbf{P3 Resolution refutation} $-$ FOPL, CNF, \textbf{negate the goal},
        resolve on complementary literals with the MGU, derive $\square$.
  \item \textbf{P4 Chaining} $-$ forward from the facts, or backward from the goal,
        on definite clauses.
  \item \textbf{P5 Bayes} $-$
        $P(A|B) = \dfrac{P(B|A)P(A)}{P(B|A)P(A) + P(B|\neg A)P(\neg A)}$, and for a
        network $P(x_1 \ldots x_n) = \prod P(x_i \mid parents(x_i))$.
\end{itemize}

{\color{sub}\footnotesize Notation used throughout: constants start with a capital,
variables are $x, y, z$, and a clause is written as a disjunction of literals. Beside
every resolvent, \mk{(i,\,j)} names the two clauses it came from and
$[t/x]$ the substitution used. That annotation is where a third of the marks sit.}

\hr

%=======================================================================
\T{P1. Truth tables: WFF, validity and tautology}

\creamq{Define a WFF with its properties. Prove that $(A \wedge (A \rightarrow B)) \rightarrow B$ is a tautology}{\m{4+4}}
{\color{sub}\footnotesize \tP{1} \yr{\textbf{\textit{71 Bh}} \m{4+4}}}

\begin{asked}{2071 Bhadra \textperiodcentered{} CT 78506 \textperiodcentered{} [4+4]}
Define a WFF with its properties. Prove that the proposition
(A $\wedge$ (A $\rightarrow$ B) $\rightarrow$ B) is a tautology.
\end{asked}

\sbsr{0.46}{%
\lead{Solution: build the table}
\begin{itemize}
  \item Two symbols, so $2^2 = 4$ rows. Work left to right, one column per sub-formula.
\end{itemize}
\vspace{1pt}
\begin{tabularx}{\linewidth}{@{}L{0.7cm}L{0.7cm}L{1.5cm}L{2.0cm}X@{}}
\toprule
\hrow \thead{A} & \thead{B} & \thead{$A \rightarrow B$} & \thead{$A \wedge (A \rightarrow B)$} & \thead{$\ldots \rightarrow B$} \\
\midrule
T & T & T & T & \textbf{T} \\
T & F & F & F & \textbf{T} \\
F & T & T & F & \textbf{T} \\
F & F & T & F & \textbf{T} \\
\bottomrule
\end{tabularx}
\vspace{2pt}
\begin{itemize}
  \item The last column is \textbf{T in every row}, so the proposition is
        \mk{valid: a tautology}. \checkmark
  \item Rows 2--4 are true \emph{vacuously}: the antecedent
        $A \wedge (A \rightarrow B)$ is false, and $F \rightarrow$ anything is T.
        \mk{Say that} or the last mark is at risk.
\end{itemize}}{%
\lead{The definition half of the answer}
\begin{itemize}
  \item A \textbf{well-formed formula} is a string built only by the recursive rules in
        \S4.2: a symbol; $\neg\alpha$; $(\alpha)$; and
        $(\alpha \wedge \beta)$, $(\alpha \vee \beta)$, $(\alpha \rightarrow \beta)$,
        $(\alpha \leftrightarrow \beta)$ for WFFs $\alpha, \beta$ $-$
        \mk{and nothing else}.
  \item \textbf{Properties to list}: it is finite; it is uniquely parsed (brackets or
        precedence fix the reading); it has a truth value under every interpretation;
        and it is exactly the set of strings the inference rules may be applied to.
  \item $A \wedge \rightarrow B$ and $(A \vee B$ are \textbf{not} WFFs: no rule builds
        them. Giving one non-example is worth a mark.
  \item \textbf{This formula is modus ponens written as a single proposition}, which is
        the point of the question: an inference rule is sound exactly when its
        implication form is a tautology. Cross-refer to the rules table and the two marks
        for ``properties'' are safe.
\end{itemize}}

%=======================================================================
\T{P2. Conversion to CNF}

{\color{sub}\footnotesize Four papers give a formula and ask for its clause form
\tF{4}. Nine more ask for the steps in the abstract; those are answered by \S4.4 plus any
one example below.}

\creamq{P $\vee$ ($\neg$P $\wedge$ Q $\wedge$ R), and ``Everyone who loves all animals is loved by someone''}{\m{2+2+2+3}}
{\color{sub}\footnotesize \yr{\textbf{\textit{73 Bh}} \m{2+2+2+3}}}

\begin{asked}{2073 Bhadra \textperiodcentered{} CT 78506 \textperiodcentered{} [2+2+2+3]}
Why is conjunctive normal form required? Explain all the steps to convert to CNF, and
transform the given sentences into CNF:
P $\vee$ ($\neg$P $\wedge$ Q $\wedge$ R); and ``Everyone who loves all animals is loved
by someone.''
\end{asked}

\sbs{%
\lead{Part 1: \; $P \vee (\neg P \wedge Q \wedge R)$}
\begin{itemize}
  \item No $\leftrightarrow$, no $\rightarrow$, no negation to move: go straight to
        \textbf{step 7}, distribute $\vee$ over $\wedge$.
  \item $P \vee (\neg P \wedge Q \wedge R)$
        \[ = (P \vee \neg P) \wedge (P \vee Q) \wedge (P \vee R) \]
  \item $P \vee \neg P$ is a \textbf{tautology}, true in every model, so the clause adds
        nothing and is dropped:
        \[ \boxed{(P \vee Q) \wedge (P \vee R)} \]
  \item \textbf{Two clauses}: $\{P, Q\}$ and $\{P, R\}$.
  \item \mk{Dropping the tautologous clause is worth doing and worth
        saying.} Leaving it in is not wrong, but the simplification is the
        part that shows the definition was understood.
\end{itemize}}{%
\lead{Part 2: the quantified sentence}
\begin{itemize}
  \item \textbf{FOPL}:
        $\forall x\,[\forall y\; Animal(y) \rightarrow Loves(x,y)]
         \rightarrow [\exists z\; Loves(z,x)]$
  \item \textbf{Eliminate $\rightarrow$} (both of them):
        $\forall x\, \neg[\forall y\; \neg Animal(y) \vee Loves(x,y)]
         \vee \exists z\; Loves(z,x)$
  \item \textbf{Move $\neg$ inwards}, dual then De Morgan:
        $\forall x\, [\exists y\; Animal(y) \wedge \neg Loves(x,y)]
         \vee \exists z\; Loves(z,x)$
  \item \textbf{Skolemize.} Both $\exists$ sit inside $\forall x$, so both become
        \mk{functions of $x$}: $y \Rightarrow F(x)$, $z \Rightarrow G(x)$.
  \item \textbf{Drop $\forall$, distribute, split}:
  \begin{itemize}
    \item $Animal(F(x)) \vee Loves(G(x), x)$
    \item $\neg Loves(x, F(x)) \vee Loves(G(x), x)$
  \end{itemize}
  \item \mk{$F$ and $G$ must be functions, not constants.} A constant
        would claim one single animal, and one single admirer, for the whole universe.
\end{itemize}}

\penalty0\vspace{2pt}

\creamq{P $\rightarrow$ ((Q $\wedge$ $\neg$R) $\rightarrow$ S)}{\m{6}}
{\color{sub}\footnotesize \yr{\textit{74 Ma} \m{6}}}

\begin{asked}{2074 Magh \textperiodcentered{} CT 78506 \textperiodcentered{} [6]}
Why is conjunctive normal form required? Explain all the steps to convert to CNF, and
transform the given sentence into CNF: \; P $\rightarrow$ ((Q $\wedge$ $\neg$R)
$\rightarrow$ S)
\end{asked}

\sbs{%
\lead{Solution}
\begin{itemize}
  \item \textbf{Inner implication first}:
        $(Q \wedge \neg R) \rightarrow S \;=\; \neg(Q \wedge \neg R) \vee S$
  \item \textbf{Outer implication}:
        $P \rightarrow [\neg(Q \wedge \neg R) \vee S]
        \;=\; \neg P \vee \neg(Q \wedge \neg R) \vee S$
  \item \textbf{De Morgan}, then double negation:
        $\neg(Q \wedge \neg R) = \neg Q \vee R$
  \item \[ \boxed{\neg P \vee \neg Q \vee R \vee S} \]
\end{itemize}}{%
\lead{The point of this one}
\begin{itemize}
  \item Nothing to distribute, so the answer is \mk{a single clause}
        with four literals. A conjunction of one clause is still CNF.
  \item Students who expect several clauses often invent a split that is not there.
        \textbf{$\vee$ distributes over $\wedge$, never the reverse}, and there is no
        $\wedge$ left at the top level here.
  \item Written as a set: $\{\neg P, \neg Q, R, S\}$. It is a
        \textbf{definite clause} $-$ two positive literals, so no; it is a plain clause.
        (Definite needs \emph{exactly one} positive literal, and this has two, $R$ and
        $S$.) A one-line check like that is what \yr{\textbf{78 Ch}}'s Horn-clause part
        is testing.
\end{itemize}}

\penalty0\vspace{2pt}

\creamq{P $\rightarrow$ ((Q $\wedge$ $\neg$R) $\leftrightarrow$ S)}{\m{8}\m{2+2+4}}
{\color{sub}\footnotesize \yr{\textbf{\textit{72 Ash}} \m{8}, 79 Jth \m{2+2+4}}
\gap$\cdot$\gap 79 Jth pairs it with
$\neg(P \wedge Q) \vee (P \vee S) \rightarrow R$, worked below}

\begin{asked}{2072 Ashwin \textperiodcentered{} CT 78506 \textperiodcentered{} [8]}
Why is conjunctive normal form required? Explain all the steps to convert to CNF, and
transform the given sentence into CNF: \; P $\rightarrow$ ((Q $\wedge$ $\neg$R)
$\leftrightarrow$ S)
\end{asked}

\sbsr{0.54}{%
\lead{Solution: the biconditional first}
\begin{itemize}
  \item Write $M = Q \wedge \neg R$ to keep the algebra readable.
  \item \textbf{Step 1, eliminate $\leftrightarrow$}:
        $M \leftrightarrow S \;=\; (M \rightarrow S) \wedge (S \rightarrow M)$
  \item \textbf{Step 2, eliminate $\rightarrow$} inside:
        $= (\neg M \vee S) \wedge (\neg S \vee M)$
  \item Substitute $M$ back and apply De Morgan to $\neg M$:
        \[ (\neg Q \vee R \vee S) \wedge (\neg S \vee (Q \wedge \neg R)) \]
  \item \textbf{Step 7, distribute} in the right half:
        $\neg S \vee (Q \wedge \neg R) = (\neg S \vee Q) \wedge (\neg S \vee \neg R)$
  \item So the inner formula is
        $(\neg Q \vee R \vee S) \wedge (\neg S \vee Q) \wedge (\neg S \vee \neg R)$.
  \item \textbf{Now the outer implication}: $P \rightarrow X = \neg P \vee X$, and
        $\neg P$ distributes into every clause of $X$:
  \[ \boxed{\begin{aligned}
     &(\neg P \vee \neg Q \vee R \vee S) \;\wedge \\
     &(\neg P \vee \neg S \vee Q) \;\wedge \\
     &(\neg P \vee \neg S \vee \neg R)
     \end{aligned}} \]
  \item \textbf{Three clauses.}
\end{itemize}}{%
\lead{\mk{Where this one goes wrong}}
\begin{itemize}
  \item \textbf{Expanding $\leftrightarrow$ the other way.} $A \leftrightarrow B$ is also
        $(A \wedge B) \vee (\neg A \wedge \neg B)$. That is correct logic but it is a
        \emph{disjunction} of conjunctions, so it takes twice the work to push back into
        CNF. \mk{Always use the two-implications form.}
  \item \textbf{Forgetting to distribute $\neg P$ into all three clauses.} $\neg P \vee
        (C_1 \wedge C_2 \wedge C_3)$ is not $\neg P \vee C_1 \wedge C_2 \wedge C_3$; the
        $\neg P$ joins \emph{each} clause.
  \item \textbf{Sanity check that costs 10 seconds}: set $P$ false. Every clause contains
        $\neg P$, so all three are satisfied $-$ which is right, because
        $P \rightarrow \ldots$ is vacuously true. Set $P$ true, $S$ true, $Q$ false: clause
        2 becomes $F \vee F \vee F$, false, and indeed $S$ true with $Q$ false breaks the
        biconditional. \mk{The CNF behaves like the original.}
\end{itemize}}

\penalty0\vspace{2pt}

\creamq{$\neg$(P $\wedge$ Q) $\vee$ (P $\vee$ S) $\rightarrow$ R}{\m{2+2+4}}
{\color{sub}\footnotesize \yr{79 Jth \m{2+2+4}} \gap$\cdot$\gap the paper's part (i);
part (ii) is the biconditional worked above}

\begin{asked}{2079 Jestha \textperiodcentered{} CT 653 \textperiodcentered{} [2+2+4]}
Why is conjunctive normal form required? Explain all the steps to convert to CNF, and
transform the given sentences into CNF:
(i) $\neg$(P $\wedge$ Q) $\vee$ (P $\vee$ S) $\rightarrow$ R \quad
(ii) P $\rightarrow$ ((Q $\wedge$ $\neg$R) $\leftrightarrow$ S)
\end{asked}

\sbs{%
\lead{Solution}
\begin{itemize}
  \item $\rightarrow$ binds \textbf{loosest}, so the sentence reads
        $\big[\neg(P \wedge Q) \vee (P \vee S)\big] \rightarrow R$.
        \mk{Reading the brackets wrongly here changes the answer};
        state the reading before starting.
  \item \textbf{Eliminate $\rightarrow$}:
        $\neg\big[\neg(P \wedge Q) \vee P \vee S\big] \vee R$
  \item \textbf{Move $\neg$ in} (De Morgan on the bracket):
        $\big[(P \wedge Q) \wedge \neg P \wedge \neg S\big] \vee R$
  \item \textbf{Distribute} $\vee R$ over the four conjuncts:
  \[ \boxed{(P \vee R) \wedge (Q \vee R) \wedge (\neg P \vee R) \wedge (\neg S \vee R)} \]
  \item \textbf{Four clauses.}
\end{itemize}}{%
\lead{Reading the answer back}
\begin{itemize}
  \item Clauses 1 and 3 are $(P \vee R)$ and $(\neg P \vee R)$. Resolving them gives $R$
        on its own, so \mk{the whole sentence is equivalent to $R$}.
  \item That is correct: the antecedent
        $\neg(P \wedge Q) \vee P \vee S$ is a \textbf{tautology} (if $P$ is false the
        first disjunct is true, if $P$ is true the second is), and
        $\text{True} \rightarrow R$ is just $R$.
  \item \mk{Saying this earns the last mark} and proves the
        conversion was not mechanical. It is also the cleanest possible demonstration of
        why resolution wants CNF: two clauses in, one simpler clause out.
\end{itemize}}

%=======================================================================
\T{P3. Translating English into FOPL}

\creamq{Convert the following into FOL}{\m{2}}
{\color{sub}\footnotesize \yr{\textbf{\textit{76 Bh}} \m{2}} \gap$\cdot$\gap
one paper, 2 marks, so no tier chip \gap$\cdot$\gap but
\mk{the same skill is step 1 of all 33 resolution proofs}, which is the
real reason to practise it}

\begin{asked}{2076 Bhadra \textperiodcentered{} CT 78506 \textperiodcentered{} [2]}
Convert the following into FOL.
\begin{enumerate}[label=\alph*), leftmargin=1.6em]
  \item Arya hates Joffery and Cersei.
  \item No Lannister loves Bran.
  \item Sansa has at least one sister.
  \item Every Stark hates at least one Lannister.
\end{enumerate}
\end{asked}

\sbs{%
\lead{Solution}
\begin{itemize}
  \item \textbf{a)} Two separate facts, joined by $\wedge$. No quantifier: both are named
        individuals.
        \[ hates(Arya, Joffery) \wedge hates(Arya, Cersei) \]
  \item \textbf{b)} ``No A loves B'' $=$ there is no A that loves B, or equivalently every
        A fails to love B:
        \[ \neg \exists x\, \big[Lannister(x) \wedge loves(x, Bran)\big] \]
        \[ \equiv\; \forall x\, \big[Lannister(x) \rightarrow \neg loves(x, Bran)\big] \]
  \item \textbf{c)} ``at least one'' $=$ $\exists$, and $\exists$ takes $\wedge$:
        \[ \exists x\, \big[sister(x) \wedge sisterOf(x, Sansa)\big] \]
  \item \textbf{d)} \mk{$\forall$ outside, $\exists$ inside}, because
        each Stark may hate a different Lannister:
        \[ \forall x\, \Big[Stark(x) \rightarrow \exists y\,
           \big(Lannister(y) \wedge hates(x,y)\big)\Big] \]
\end{itemize}}{%
\lead{\mk{The four rules this question tests}}
\begin{itemize}
  \item $\forall$ takes $\rightarrow$; $\exists$ takes $\wedge$. Every wrong answer to
        (c) or (d) breaks this.
  \item \textbf{Quantifier order carries meaning.}
        $\forall x \exists y$ (each Stark hates \emph{some} Lannister) is not
        $\exists y \forall x$ (one Lannister is hated by \emph{all} Starks). The English
        ``at least one'' after ``every'' fixes the order as $\forall\exists$.
  \item \textbf{``No'' negates the whole existential}, not the predicate. Writing
        $\forall x\, Lannister(x) \wedge \neg loves(x,Bran)$ says everything in the
        universe is a Lannister.
  \item Both forms of (b) are correct and \textbf{equivalent by the dual}
        $\neg\exists \equiv \forall\neg$. Giving both, with that one-line justification,
        is a complete answer to a 2-mark part.
  \item \textbf{Predicate names are your choice}, but stay consistent: $sisterOf(x,y)$
        must mean the same thing every time it appears.
\end{itemize}}

%=======================================================================
\T{P4. Unification: finding the MGU}

{\color{sub}\footnotesize Never asked on its own, but every resolution step needs one and
the marks come off when it is wrong. Four standard pairs, in the form the papers and the
textbook use.}

\sbs{%
\lead{Worked}
\begin{itemize}
  \item \textbf{1.} $King(x)$, $King(John)$ \; $\rightarrow$ \;
        $\theta = [John/x]$.
  \item \textbf{2.} $P(x, y)$, $P(a, f(z))$ \; $\rightarrow$ \;
        $\theta = [a/x,\; f(z)/y]$. \; $z$ stays free: binding it would be
        \emph{less} general.
  \item \textbf{3.} $P(x, f(y))$, $P(a, f(g(x)))$ \; $\rightarrow$ \;
        start with $[a/x]$, then $y$ must match $g(x)$, and $x$ is already $a$:
        \[ \theta = [a/x,\; g(a)/y] \]
  \item \textbf{4.} $Q(a, g(x,a), f(y))$, $Q(a, g(f(b),a), x)$ \; $\rightarrow$ \;
        $x$ must be $f(b)$ from the second argument, and then $f(y)$ must match
        $x = f(b)$, giving $y = b$:
        \[ \theta = [f(b)/x,\; b/y] \]
\end{itemize}}{%
\lead{\mk{The three ways it fails}}
\begin{itemize}
  \item \textbf{Different predicate or arity.} $P(x)$ and $Q(x)$ never unify; neither do
        $P(x)$ and $P(x,y)$.
  \item \textbf{Two different constants.} $King(John)$ and $King(Ram)$ fail: a constant
        cannot be substituted.
  \item \textbf{Occurs check.} $x$ cannot bind to $f(x)$; the substitution would never
        terminate.
  \item \mk{Apply each binding to everything already found}, as in
        example 3. Reporting $[a/x, g(x)/y]$ instead of $[a/x, g(a)/y]$ is the standard
        half-mark loss.
  \item \textbf{Standardise apart before unifying two clauses.} If both clauses call their
        variable $x$, rename one to $x'$ first, or the unifier will tie together two
        variables that have nothing to do with each other.
\end{itemize}}

%=======================================================================
\T{P5. Resolution refutation \gap\textcolor{sub}{Family A: likes all kinds of food}}

{\color{sub}\footnotesize \mk{7 papers, 5 wordings, 55 marks} $-$ the
single most-repeated premise set in the subject \yr{\textbf{72 Ash} \m{7},
\textbf{74 Bh} \m{8}, \textbf{\textit{73 Bh}} \m{8}, \textbf{\textit{71 Bh}} \m{8},
\textbf{\texttt{80 Bh}} \m{8}, \texttt{80 Ba} \m{8}, \textit{72 Ma} \m{6+4}}. Work it once
properly and five papers are answered.}

\begin{asked}{2072 Ashwin \textperiodcentered{} Q10 \textperiodcentered{} [7] \; also 2074 Bhadra [8], 2073 Bhadra CT 78506 [8]}
Assume the following facts: John likes all kinds of food; apples are food; chicken is
food; anything anyone eats and isn't killed by is food; Bill eats peanuts and is still
alive; Sue eats everything Bill eats. Prove that John likes peanuts using resolution
refutation.
\end{asked}

\sbsr{0.47}{%
\lead{Step 1 \; Facts into FOPL}
\begin{itemize}
  \item[1.] $\forall x\; food(x) \rightarrow likes(John, x)$
  \item[2.] $food(Apple)$
  \item[3.] $food(Chicken)$
  \item[4.] $\forall y \forall x\; \big[eats(y,x) \wedge \neg killed(y)\big]
            \rightarrow food(x)$
  \item[5.] $eats(Bill, Peanuts) \wedge \neg killed(Bill)$
  \item[6.] $\forall x\; eats(Bill, x) \rightarrow eats(Sue, x)$
\end{itemize}
{\color{sub}\footnotesize ``is still alive'' is written $\neg killed$, not a new predicate
$alive$. One predicate fewer, one resolution step fewer.}
\lead{Step 2 \; Into CNF}
\begin{itemize}
  \item[1.] $\neg food(x) \vee likes(John, x)$
  \item[2.] $food(Apple)$
  \item[3.] $food(Chicken)$
  \item[4.] $\neg eats(y,x) \vee killed(y) \vee food(x)$
  \item[5a.] $eats(Bill, Peanuts)$
  \item[5b.] $\neg killed(Bill)$
  \item[6.] $\neg eats(Bill, x) \vee eats(Sue, x)$
\end{itemize}
\lead{Step 3 \; Negate the goal}
\begin{itemize}
  \item[7.] $\neg likes(John, Peanuts)$
\end{itemize}}{%
\lead{Step 4 \; Resolution}
\begin{itemize}
  \item[8.] $\neg food(Peanuts)$ \hfill \yr{(7,\,1) $[Peanuts/x]$}
  \item[9.] $\neg eats(y, Peanuts) \vee killed(y)$ \hfill \yr{(8,\,4) $[Peanuts/x]$}
  \item[10.] $killed(Bill)$ \hfill \yr{(9,\,5a) $[Bill/y]$}
  \item[11.] $\square$ \hfill \yr{(10,\,5b)}
\end{itemize}
\lead{Final answer}
\begin{itemize}
  \item The empty clause is derived, so $KB \wedge \neg goal$ is unsatisfiable.
        \mk{``John likes peanuts'' is proved.} \checkmark
  \item Clauses \textbf{2, 3 and 6 are never used}: apples, chicken and Sue are
        distractors. Saying so is a mark in a 8-mark answer, not a risk.
\end{itemize}
\lead{\mk{The step the textbook gets wrong}}
\begin{itemize}
  \item Fact 4 has the shape $\neg[A \wedge \neg B] \vee C$, which is
        \mk{$\neg A \vee B \vee C$: one clause}.
  \item \emph{Insights} splits it into two, $\neg eats(y,x) \vee food(x)$ and
        $\neg killed(y) \vee food(x)$. That is not a valid equivalence $-$ it asserts that
        merely \emph{not being killed} makes something food $-$ and the book's own
        resolution graph then quietly uses the correct single clause.
  \item Write the one clause. If the marker expects the book's two, the proof still
        closes, and yours is the sound one.
\end{itemize}}

\penalty0\vspace{2pt}

\creamq{The four renamings}{\m{8}\m{6+4}}
{\color{sub}\footnotesize same proof, different names \gap$\cdot$\gap
\mk{the last one is not the same} and is the only place in the family where
the printed goal does not follow}

\begin{asked}{2071 Bhadra \textperiodcentered{} CT 78506 \textperiodcentered{} [8]}
Assume the following facts: a) Ram likes all kinds of food; b) Apples are food;
c) Chicken is food; d) Anything anyone eats and is not killed by is food; e) Krishna eats
peanuts and is still alive; f) Radha eats everything Krishna eats. Prove the statement
``Ram likes peanut'' using resolution.
\end{asked}

\begin{asked}{2080 Bhadra \textperiodcentered{} CT 710 \textperiodcentered{} [8]}
Assume the following facts: i) Ram likes all kinds of food; ii) apples and vegetables are
food; iii) anything anyone eats and not killed is food; iv) Anil eats peanuts and is still
alive; v) Ram eats everything that Anil eats. Prove that ``Ram likes peanuts'' using
resolution refutation.
\end{asked}

\begin{asked}{2080 Baishakh \textperiodcentered{} CT 710 \textperiodcentered{} [8]}
Assume the following facts: a) Sita likes all kinds of food; b) Momo and Burger are food;
c) anything anyone eats and isn't killed by is food; d) Bill eats mushroom and is still
alive; e) Sue eats everything Bill eats. Prove that ``Sita likes mushroom'' using
resolution refutation method.
\end{asked}

\sbs{%
\lead{What changes, and what does not}
\vspace{-1pt}
\begin{tabularx}{\linewidth}{@{}L{1.75cm}L{2.2cm}L{2.0cm}X@{}}
\toprule
\hrow \thead{Paper} & \thead{Likes all food} & \thead{Eats it} & \thead{Goal} \\
\midrule
\textbf{72 Ash}, \textbf{74 Bh}, \textbf{\textit{73 Bh}} & John & Bill (Sue copies) & peanuts \\
\textbf{\textit{71 Bh}} & Ram & Krishna (Radha copies) & peanuts \\
\textbf{\texttt{80 Bh}} & Ram & Anil (\textbf{Ram} copies) & peanuts \\
\texttt{80 Ba} & Sita & Bill (Sue copies) & mushroom \\
\bottomrule
\end{tabularx}
\vspace{2pt}
\begin{itemize}
  \item \mk{The proof is identical in all four}: negate, resolve with
        clause 1 to get $\neg food(g)$, resolve with clause 4 to get
        $\neg eats(y,g) \vee killed(y)$, then the two halves of the eating fact.
        Substitute the names and the four steps are unchanged.
  \item \yr{\textbf{\texttt{80 Bh}}} is the one worth a second look: there the copier is
        \textbf{Ram himself}, the same person who likes all food. It changes nothing $-$
        the copying clause is still unused $-$ but it looks like it should, and that is
        the trap.
  \item \yr{\texttt{80 Ba}} swaps apples and chicken for momo and burger. Also unused.
\end{itemize}}{%
\lead{\mk{2072 Magh (CT 78506) is a different problem}}
\begin{itemize}
  \item It adds a fifth premise, \emph{``if a person likes a food it means that person has
        eaten it''}, and asks for \textbf{Shailendri} $-$ the copier $-$ to
        \emph{\textbf{like}} chicken.
  \item Take the premises literally:
  \begin{itemize}
    \item $food(Chicken)$ and $\forall x\, food(x) \rightarrow likes(Bhogendra, x)$
          give $likes(Bhogendra, Chicken)$;
    \item $\forall x \forall y\, likes(x,y) \rightarrow eats(x,y)$
          gives $eats(Bhogendra, Chicken)$;
    \item $\forall x\, eats(Bhogendra,x) \rightarrow eats(Shailendri,x)$
          gives \mk{$eats(Shailendri, Chicken)$}.
  \end{itemize}
  \item And there it stops. \textbf{Nothing in the premise set makes Shailendri
        \emph{like} anything}: only Bhogendra has a likes-rule, and the new premise runs
        \emph{likes $\rightarrow$ eats}, not back.
  \item \textbf{How to answer it}: prove $eats(Shailendri, Chicken)$ in full, then add one
        line $-$ \emph{``to reach \emph{likes} the fifth premise must be read as a
        biconditional, $likes(x,y) \leftrightarrow eats(x,y)$; with that reading one more
        resolution step gives $likes(Shailendri, Chicken)$.''}
  \item \mk{Do not silently reverse the implication and hope.} Stating
        the gap is worth more than a proof that assumes it away.
\end{itemize}}

\penalty0\vspace{2pt}

\begin{asked}{2072 Magh \textperiodcentered{} CT 78506 \textperiodcentered{} [6+4]}
Bhogendra likes all kinds of food. Oranges are food. Chicken is food. Anything anyone eats
and isn't killed by is food. If a person likes a food it means that person has eaten it.
Jogendra eats peanuts and is still alive. Shailendri eats everything Bhogendra eats.
Express them in FOPL, put them into clause form, and use resolution to prove that
Shailendri likes chicken.
\end{asked}

%=======================================================================
\T{P5. Resolution refutation \gap\textcolor{sub}{Family B: Charlie is a horse}}

{\color{sub}\footnotesize \mk{6 papers, 3 wordings, 46 marks}
\yr{\textbf{73 Bh} \m{8}, 78 Po \m{8}, \textbf{\textit{74 Bh}} \m{8}, 79 Jth \m{6},
\textbf{\texttt{79 Bh}} \m{8}, 78 Ba \m{8}}. The set that tests
\textbf{inverse relations} and, quietly, whether you can spot an unused premise.}

\begin{asked}{2073 Bhadra \textperiodcentered{} Q10 \textperiodcentered{} [8] \; also 2078 Poush [8], 2074 Bhadra CT 78506 [8], 2079 Jestha [6]}
Assume the following facts: i) horses, cows, pigs are mammals; ii) an offspring of a horse
is a horse; iii) Bluebeard is a horse; iv) Bluebeard is Charlie's parent; v) offspring and
parent are inverse relations; vi) every mammal has a parent. Prove Charlie is a horse
using resolution refutation.
\end{asked}

\sbsr{0.50}{%
\lead{Step 1 \; Facts into FOPL}
\begin{itemize}
  \item[1.] $\forall x\; horse(x) \rightarrow mammal(x)$, and the same for $cow$, $pig$
  \item[2.] $\forall x \forall y\; \big[offspring(x,y) \wedge horse(y)\big]
            \rightarrow horse(x)$
  \item[3.] $horse(Bluebeard)$
  \item[4.] $parent(Bluebeard, Charlie)$
  \item[5.] $\forall x \forall y\; offspring(x,y) \leftrightarrow parent(y,x)$
  \item[6.] $\forall x\; mammal(x) \rightarrow \exists y\; parent(y, x)$
\end{itemize}
\lead{Step 2 \; Into CNF}
\begin{itemize}
  \item[1a.] $\neg horse(x) \vee mammal(x)$ \; (1b, 1c likewise for cow, pig)
  \item[2.] $\neg offspring(x,y) \vee \neg horse(y) \vee horse(x)$
  \item[3.] $horse(Bluebeard)$
  \item[4.] $parent(Bluebeard, Charlie)$
  \item[5a.] $\neg offspring(x,y) \vee parent(y,x)$
  \item[5b.] $\neg parent(y,x) \vee offspring(x,y)$
  \item[6.] $\neg mammal(x) \vee parent\big(f(x),\, x\big)$
\end{itemize}
{\color{sub}\footnotesize The $\exists y$ in premise 6 sits inside $\forall x$, so it
skolemizes to a \textbf{function} $f(x)$: each mammal has its own parent.
\mk{\emph{Insights} writes a free variable $y$ here}, which claims every
object is a parent of every mammal. It changes nothing in this proof, because clause 6 is
never used, but it is wrong and in another problem it would prove anything.}}{%
\lead{Step 3 \; Negate the goal}
\begin{itemize}
  \item[7.] $\neg horse(Charlie)$
\end{itemize}
\lead{Step 4 \; Resolution}
\begin{itemize}
  \item[8.] $\neg offspring(Charlie, y) \vee \neg horse(y)$
        \hfill \yr{(7,\,2) $[Charlie/x]$}
  \item[9.] $\neg offspring(Charlie, Bluebeard)$
        \hfill \yr{(8,\,3) $[Bluebeard/y]$}
  \item[10.] $\neg parent(Bluebeard, Charlie)$
        \hfill \yr{(9,\,5b) $[Charlie/x, Bluebeard/y]$}
  \item[11.] $\square$ \hfill \yr{(10,\,4)}
\end{itemize}
\lead{Final answer}
\begin{itemize}
  \item \mk{``Charlie is a horse'' is proved.} \checkmark \; Four
        resolution steps.
  \item \textbf{Premises 1 and 6 are never used.} The mammal chain is a complete red
        herring for this goal: being a horse follows from the offspring rule alone.
        \mk{Say it in one line} $-$ examiners set the extra premises to
        see whether the proof is understood or remembered.
  \item \textbf{The whole problem is premise 5.} Written as a biconditional it yields the
        two clauses 5a and 5b, and only \textbf{5b} is used: it turns the given
        $parent(Bluebeard, Charlie)$ into the $offspring(Charlie, Bluebeard)$ that rule 2
        needs. A one-directional reading of ``inverse relations'' cannot close the proof.
  \item \textbf{Watch the argument order}: $offspring(x,y)$ means \emph{$x$ is the
        offspring of $y$}, so its inverse is $parent(y,x)$, not $parent(x,y)$. Getting
        this backwards is the commonest failure on this set.
\end{itemize}}

\penalty0\vspace{2pt}

\creamq{The two renamings}{\m{8}}
{\color{sub}\footnotesize \yr{\textbf{\texttt{79 Bh}} \m{8}} asks for a \emph{mammal},
\yr{78 Ba \m{8}} renames every noun}

\begin{asked}{2079 Bhadra \textperiodcentered{} CT 710 \textperiodcentered{} [8]}
Prove that ``Charlie is a mammal'' with given propositions: a) cows, pigs and horses are
mammal; b) the child of horse is a horse; c) Bluebeard is a horse; d) Bluebeard is
Charlie's father; e) child and father are inverse relation; f) every mammal has a father.
\end{asked}

\begin{asked}{2078 Baishakh \textperiodcentered{} [8]}
Assume the following facts: (i) Cow, Buffalo, Bull are mammals; (ii) an offspring of a Cow
is Cow; (iii) Kali is a Cow; (iv) Kali is Taarey's parent; (v) offspring and parent are
inverse relations; (vi) every mammal has a parent. By using resolution refutation, prove
that Taarey is a Cow.
\end{asked}

\sbs{%
\lead{2079 Bhadra: one extra step, and premise (a) stops being a distractor}
\begin{itemize}
  \item The goal is now $mammal(Charlie)$, so negate to
        \textbf{7.} $\neg mammal(Charlie)$.
  \item[] \textbf{8.} $\neg horse(Charlie)$ \hfill \yr{(7,\,1a) $[Charlie/x]$}
  \item[] and from there the proof above runs unchanged, steps 8--11 becoming 9--12.
  \item \mk{Five steps instead of four.} The mammal rule that was
        unused in the ``horse'' version is now the first step of the ``mammal'' one $-$
        the same six premises, one different goal, a different subset used.
  \item \emph{child} and \emph{father} replace \emph{offspring} and \emph{parent}. Same
        two clauses from the inverse relation, same argument-order trap.
\end{itemize}}{%
\lead{2078 Baishakh: a pure renaming}
\vspace{-1pt}
\begin{tabularx}{\linewidth}{@{}L{2.6cm}XX@{}}
\toprule
\hrow & \thead{2073 Bhadra} & \thead{2078 Baishakh} \\
\midrule
Species & horse & Cow \\
Also mammals & cows, pigs & Buffalo, Bull \\
The parent & Bluebeard & Kali \\
The offspring & Charlie & Taarey \\
Goal & $horse(Charlie)$ & $Cow(Taarey)$ \\
\bottomrule
\end{tabularx}
\vspace{2pt}
\begin{itemize}
  \item Substitute and the four resolution steps are character-for-character the same.
  \item \textbf{Buffalo and Bull are the distractors} this time, and premise (vi) is
        again unused.
  \item \mk{Answer to give}: $\square$ derived in 4 steps, so
        ``Taarey is a Cow'' is proved. \checkmark
\end{itemize}}

%=======================================================================
\T{P5. Resolution refutation \gap\textcolor{sub}{Family C: selling weapons to a hostile nation}}

{\color{sub}\footnotesize \mk{7 papers, 6 wordings, 55 marks}
\yr{\textbf{69 Bh} \m{10}, 75 Ba \m{8}, \textbf{77 Ch} \m{8}, 82 Ka \m{8},
81 Ash \m{2+6}, \textbf{\texttt{81 Bh}} \m{9}, \textbf{\texttt{82 Bh}} \m{6}}. The set
that tests \textbf{existential instantiation}, and the one paper
\yr{(81 Ash)} that asks for \textbf{forward chaining} instead of refutation.}

\begin{asked}{2077 Chaitra \textperiodcentered{} Q10 \textperiodcentered{} [8] \; also 2082 Kartik [8]}
It is a crime for an American to sell weapons to hostile nations. Nono has some missiles.
All the missiles owned by Nono were sold to it by Colonel West. Missiles are weapons. An
enemy of America counts as hostile. Colonel West is an American. The country Nono is an
enemy of America. Prove that Colonel West is a criminal by using FOPL based Resolution
Refutation.
\end{asked}

\sbsr{0.49}{%
\lead{Step 1 \; Facts into FOPL}
\begin{itemize}
  \item[1.] $\forall p \forall q \forall r\; \big[American(p) \wedge Weapon(q) \wedge
            Sells(p,q,r) \wedge Hostile(r)\big] \rightarrow Criminal(p)$
  \item[2.] $\exists x\; Missile(x) \wedge Owns(Nono, x)$
  \item[3.] $\forall x\; \big[Missile(x) \wedge Owns(Nono,x)\big]
            \rightarrow Sells(West, x, Nono)$
  \item[4.] $\forall x\; Missile(x) \rightarrow Weapon(x)$
  \item[5.] $\forall x\; Enemy(x, America) \rightarrow Hostile(x)$
  \item[6.] $American(West)$
  \item[7.] $Enemy(Nono, America)$
\end{itemize}
\lead{Step 2 \; Into CNF}
\begin{itemize}
  \item[1.] $\neg American(p) \vee \neg Weapon(q) \vee \neg Sells(p,q,r) \vee
            \neg Hostile(r) \vee Criminal(p)$
  \item[2a.] $Missile(M_1)$ \hfill \yr{skolem constant}
  \item[2b.] $Owns(Nono, M_1)$
  \item[3.] $\neg Missile(x) \vee \neg Owns(Nono,x) \vee Sells(West,x,Nono)$
  \item[4.] $\neg Missile(x) \vee Weapon(x)$
  \item[5.] $\neg Enemy(x, America) \vee Hostile(x)$
  \item[6.] $American(West)$ \qquad \textbf{7.} $Enemy(Nono, America)$
\end{itemize}
{\color{sub}\footnotesize \mk{``Nono has some missiles'' is existential.}
It is not $\forall x\, Missile(x)$, which would make everything a missile, and it is not a
named missile either $-$ the paper never names one. Skolemize to a brand-new constant
$M_1$ and it splits into two ground facts.}}{%
\lead{Step 3 \; Negate the goal}
\begin{itemize}
  \item[8.] $\neg Criminal(West)$
\end{itemize}
\lead{Step 4 \; Resolution}
\begin{itemize}
  \item[9.] $\neg American(West) \vee \neg Weapon(q) \vee \neg Sells(West,q,r) \vee
            \neg Hostile(r)$ \hfill \yr{(8,\,1) $[West/p]$}
  \item[10.] $\neg Weapon(q) \vee \neg Sells(West,q,r) \vee \neg Hostile(r)$
            \hfill \yr{(9,\,6)}
  \item[11.] $\neg Missile(q) \vee \neg Sells(West,q,r) \vee \neg Hostile(r)$
            \hfill \yr{(10,\,4) $[q/x]$}
  \item[12.] $\neg Missile(q) \vee \neg Owns(Nono,q) \vee \neg Hostile(Nono)$
            \hfill \yr{(11,\,3) $[q/x,\, Nono/r]$}
  \item[13.] $\neg Owns(Nono, M_1) \vee \neg Hostile(Nono)$
            \hfill \yr{(12,\,2a) $[M_1/q]$}
  \item[14.] $\neg Hostile(Nono)$ \hfill \yr{(13,\,2b)}
  \item[15.] $\neg Enemy(Nono, America)$ \hfill \yr{(14,\,5) $[Nono/x]$}
  \item[16.] $\square$ \hfill \yr{(15,\,7)}
\end{itemize}
\lead{Final answer}
\begin{itemize}
  \item \mk{``Colonel West is a criminal'' is proved.} \checkmark \;
        Eight resolution steps, and \textbf{every premise is used} $-$ unusual, and worth
        noticing, because it means no step can be skipped.
  \item \textbf{Work down the negated goal.} Clause 1 has five literals; resolving it
        against clause 8 first turns the proof into a checklist of four conditions, each
        discharged by one premise. Starting anywhere else wanders.
  \item In step 12 the literal $\neg Missile(q)$ appears twice after the substitution and
        \mk{collapses to one}. Duplicate literals in a clause always do.
\end{itemize}}

\penalty0\vspace{2pt}

\creamq{2069 Bhadra and 2075 Baishakh: two premises are missing}{\m{10}\m{8}}
{\color{sub}\footnotesize \yr{\textbf{69 Bh} \m{10}, 75 Ba \m{8}} \gap$\cdot$\gap
\mk{neither paper says missiles are weapons, or that an enemy is hostile}
\gap$\cdot$\gap the proof cannot close without them, so they must be stated as assumptions}

\begin{asked}{2069 Bhadra \textperiodcentered{} Q10 \textperiodcentered{} [10]}
Every American who sells weapons to hostile nations is a criminal. The country Abc is
enemy of America. All of the missiles in Abc were sold by John. John is an American. Prove
John is a criminal.
\end{asked}

\begin{asked}{2075 Baishakh \textperiodcentered{} [8]}
Every American who sells weapons to hostile nations is a criminal. The country XYZ is
enemy of America. All of its missiles in XYZ were sold by Donald, who is an American.
Prove that Donald is a criminal using FOPL based resolution refutation.
\end{asked}

\sbs{%
\lead{What is missing, and what to write}
\begin{itemize}
  \item Clause 1 needs \textbf{four} things about $p$: American, a weapon, a sale, and a
        hostile buyer. These two papers supply \emph{American}, the \emph{sale} and the
        \emph{enemy} $-$ but nothing links
  \begin{itemize}
    \item \emph{missile} to \emph{weapon}, and
    \item \emph{enemy of America} to \emph{hostile}.
  \end{itemize}
  \item \mk{Add both, and say you are adding them}:
  \begin{itemize}
    \item[] $\forall x\; Missile(x) \rightarrow Weapon(x)$
    \item[] $\forall x\; Enemy(x, America) \rightarrow Hostile(x)$
  \end{itemize}
  \item Then the proof is exactly the 8 steps above with $[John/West,\, Abc/Nono]$ or
        $[Donald/West,\, XYZ/Nono]$.
  \item \textbf{This is not pedantry.} A resolution proof is a formal derivation: if the
        bridging axiom is not in the clause set, the empty clause cannot be derived, and
        a proof that pretends otherwise is the thing being marked.
\end{itemize}}{%
\lead{\mk{Also read ``all of the missiles in Abc'' carefully}}
\begin{itemize}
  \item It says \emph{every missile in Abc was sold by John}, i.e.\
        $\forall x\, \big[Missile(x) \wedge In(x, Abc)\big] \rightarrow Sells(John, x, Abc)$.
  \item It does \textbf{not} say Abc has any. Strictly, with no missile in existence the
        antecedent is never satisfied and nothing follows $-$ so the existential
        \emph{``Abc has some missiles''} has to be assumed too, exactly as \yr{\textbf{77 Ch}}
        states it outright.
  \item \mk{Three assumptions, one line each}, then the standard proof.
        An 10-mark answer has room for them and the examiner is looking for whether the
        gap was noticed.
  \item \yr{75 Ba}'s ``All of its missiles in XYZ were sold by Donald'' is the same
        sentence with a typo; treat it identically.
\end{itemize}}

\penalty0\vspace{2pt}

\creamq{2081 Bhadra and 2082 Bhadra: the two renamings that add a premise}{\m{9}\m{6}}
{\color{sub}\footnotesize \yr{\textbf{\texttt{81 Bh}} \m{9}, \textbf{\texttt{82 Bh}} \m{6}}}

\begin{asked}{2081 Bhadra \textperiodcentered{} CT 710 \textperiodcentered{} [9]}
Every American who sells weapons to hostile nations is a criminal. Every enemy of America
is a hostile. Nono has some missiles. All missiles of Nono were sold by Gorge. Gorge is an
American. Nono is a Country. Nono is the enemy of America. Missiles are weapons. Prove
that ``Is Gorge a criminal?'' using Resolution by Refutation.
\end{asked}

\begin{asked}{2082 Bhadra \textperiodcentered{} CT 710 \textperiodcentered{} [6]}
It is a crime for a Nepali to sell weapons to an enemy of Nepal. Mr.\ A is an enemy of
Nepal. Mr.\ A has some missiles. Missiles are weapons. Mr.\ R is a Nepali. If Mr.\ A has
missiles then those missiles were sold to Mr.\ A by Mr.\ R. Prove ``Mr.\ A is a
criminal''.
\end{asked}

\sbs{%
\lead{2081 Bhadra: complete, with one spare premise}
\begin{itemize}
  \item All eight premises of the canonical set are present, plus
        \emph{``Nono is a Country''}, $Country(Nono)$, which
        \mk{appears in no rule and is never used}.
  \item Substitute $[Gorge/West]$ and the eight steps run unchanged. Answer:
        \mk{Gorge is a criminal.} \checkmark
  \item The question is phrased as \emph{``Is Gorge a criminal?''} but asks for resolution
        \emph{by refutation}, so still negate: assume $\neg Criminal(Gorge)$ and derive
        $\square$. \textbf{Do not switch to extraction} just because the sentence ends in
        a question mark.
\end{itemize}}{%
\lead{2082 Bhadra: \mk{the goal is the wrong person}}
\begin{itemize}
  \item Read the roles: \textbf{Mr.\,R} is the Nepali who \emph{sells}; \textbf{Mr.\,A}
        is the enemy who \emph{buys}. The crime rule makes the \emph{seller} the criminal.
  \item So the premises prove $Criminal(Mr.R)$, and the printed goal
        \emph{``Mr.\,A is a criminal''} does not follow.
  \item \mk{How to answer}: prove $Criminal(Mr.R)$ in full by the
        standard eight steps, then one line $-$ \emph{``the premises make the seller the
        criminal, so what follows is $Criminal(Mr.R)$; $Criminal(Mr.A)$ is not derivable
        from this set.''}
  \item The sixth premise is a conditional, \emph{if Mr.\,A has missiles then they were
        sold to him by Mr.\,R}, which is the same clause 3 as ever:
        $\neg Missile(x) \vee \neg Owns(A,x) \vee Sells(R,x,A)$.
  \item \emph{Enemy of Nepal} is used directly as the hostile predicate here, so no
        bridging axiom is needed.
\end{itemize}}

\penalty0\vspace{2pt}

%=======================================================================
\T{P6. The same premises by forward and backward chaining}

\creamq{Prove Col.\ West is criminal using forward chaining}{\m{2+6}}
{\color{sub}\footnotesize \tP{1} \yr{81 Ash \m{2+6}} \gap$\cdot$\gap
\mk{the only paper that asks for a chaining proof rather than resolution},
and 11 more ask for the forward/backward comparison this answers by example}

\begin{asked}{2081 Ashwin \textperiodcentered{} [2+6]}
The law says that it is a crime for an American to sell weapons to hostile nations. The
country Nono, an enemy of America, has some missiles, and all of its missiles were sold to
it by Colonel West, who is American. Prove that Col.\ West is criminal using
\textbf{forward chaining}.
\end{asked}

\sbsr{0.52}{%
\lead{Forward chaining, iteration by iteration}
\begin{itemize}
  \item \textbf{Start with the ground facts} (the clauses with no premises):
  \begin{itemize}
    \item[] $American(West)$, \; $Missile(M_1)$, \; $Owns(Nono, M_1)$, \;
            $Enemy(Nono, America)$
  \end{itemize}
  \item \textbf{Iteration 1} $-$ fire every rule whose premises are all present:
  \begin{itemize}
    \item rule 4 with $[M_1/x]$ \; $\rightarrow$ \; \textbf{$Weapon(M_1)$}
    \item rule 3 with $[M_1/x]$, needing $Missile(M_1)$ \emph{and} $Owns(Nono,M_1)$,
          both present \; $\rightarrow$ \; \textbf{$Sells(West, M_1, Nono)$}
    \item rule 5 with $[Nono/x]$ \; $\rightarrow$ \; \textbf{$Hostile(Nono)$}
    \item rule 1 does \textbf{not} fire yet: its four premises are not all known.
  \end{itemize}
  \item \textbf{Iteration 2} $-$ rule 1 now matches with
        $[West/p,\, M_1/q,\, Nono/r]$, because $American(West)$, $Weapon(M_1)$,
        $Sells(West,M_1,Nono)$ and $Hostile(Nono)$ are all in the fact base:
  \begin{itemize}
    \item[] $\rightarrow$ \; \mk{$Criminal(West)$}
  \end{itemize}
  \item \textbf{Goal reached, stop.} \checkmark
\end{itemize}
{\color{sub}\footnotesize Draw it as three rows: the four given facts along the bottom,
the three iteration-1 facts above them, $Criminal(West)$ at the top, with an arrow from
each premise into the fact it helped derive. That picture is the answer; see \S4.6.}}{%
\lead{The same thing backwards, for the 2 marks of comparison}
\begin{itemize}
  \item \textbf{Goal} $Criminal(West)$. Only rule 1 concludes $Criminal$, so unify
        $[West/p]$ and take its four premises as sub-goals:
  \begin{itemize}
    \item $American(West)$ \; \checkmark a fact, done;
    \item $Weapon(q)$ \; $\rightarrow$ sub-goal $Missile(q)$ \; \checkmark
          matches the fact $Missile(M_1)$, binding $[M_1/q]$;
    \item $Sells(West, M_1, r)$ \; $\rightarrow$ sub-goals $Missile(M_1)$ \checkmark
          and $Owns(Nono, M_1)$ \checkmark, binding $[Nono/r]$;
    \item $Hostile(Nono)$ \; $\rightarrow$ sub-goal $Enemy(Nono, America)$ \; \checkmark.
  \end{itemize}
  \item Every leaf is a fact, so the goal is proved.
        \mk{Four sub-goals, depth first, left to right.}
\end{itemize}
\lead{\mk{What the examiner is comparing}}
\begin{itemize}
  \item Forward derived \textbf{three facts nobody asked for} on the way
        ($Weapon$, $Sells$, $Hostile$ are all needed here, but on a bigger KB most would
        not be).
  \item Backward touched \textbf{only the four conjuncts of one rule}, and the binding
        $[M_1/q]$ found in the second sub-goal was
        \mk{carried into the third}, which is why $Sells(West,M_1,r)$
        could be matched at all.
  \item Same rules, same modus ponens, opposite ends. One sentence, and the comparison
        marks are secured.
\end{itemize}}

%=======================================================================
\T{P5. Resolution refutation \gap\textcolor{sub}{Family D: the light sleeper}}

{\color{sub}\footnotesize \mk{2 papers, 17 marks}
\yr{76 Ba \m{2+8}, \textbf{\textit{76 Bh}} \m{7}}. \textbf{The hardest set in the
chapter}, because the goal is itself an implication and the premises hide two nested
quantifiers.}

\begin{asked}{2076 Baishakh \textperiodcentered{} [2+8]}
List down the rule for Inference. Consider the following axioms: All hounds howl at night.
Anyone who has any cats will not have any mice. Light sleepers do not have anything which
howls at night. John has either a cat or a hound. Prove: ``If John is a light sleeper, then
John does not have any mice'' by resolution refutation.
\end{asked}

\sbsr{0.50}{%
\lead{Step 1 \; Facts into FOPL}
\begin{itemize}
  \item[1.] $\forall x\; Hound(x) \rightarrow Howl(x)$
  \item[2.] $\forall x \forall y\; \big[Have(x,y) \wedge Cat(y)\big] \rightarrow
            \forall z\, \big[Have(x,z) \rightarrow \neg Mouse(z)\big]$
  \item[3.] $\forall x\; LS(x) \rightarrow \forall y\,
            \big[Have(x,y) \rightarrow \neg Howl(y)\big]$
  \item[4.] $\exists x\; Have(John, x) \wedge \big[Cat(x) \vee Hound(x)\big]$
  \item[] \textbf{Goal} \;
        $LS(John) \rightarrow \forall z\,
        \big[Have(John,z) \rightarrow \neg Mouse(z)\big]$
\end{itemize}
\lead{Step 2 \; Into CNF}
\begin{itemize}
  \item[1.] $\neg Hound(x) \vee Howl(x)$
  \item[2.] $\neg Have(x,y) \vee \neg Cat(y) \vee \neg Have(x,z) \vee \neg Mouse(z)$
  \item[3.] $\neg LS(x) \vee \neg Have(x,y) \vee \neg Howl(y)$
  \item[4a.] $Have(John, a)$ \hfill \yr{skolem constant $a$}
  \item[4b.] $Cat(a) \vee Hound(a)$
\end{itemize}
\lead{Step 3 \; \mk{Negate the goal, and get three clauses}}
\begin{itemize}
  \item $\neg\big[LS(John) \rightarrow \forall z (Have(John,z) \rightarrow \neg
        Mouse(z))\big]$
  \item $= LS(John) \wedge \exists z\, \big[Have(John,z) \wedge Mouse(z)\big]$
  \item[5a.] $LS(John)$ \qquad \textbf{5b.} $Have(John, b)$ \qquad
             \textbf{5c.} $Mouse(b)$ \hfill \yr{skolem $b$}
\end{itemize}
{\color{sub}\footnotesize $\neg(P \rightarrow Q) = P \wedge \neg Q$. \mk{This
is the whole difficulty of the problem}: students negate an implication to
another implication, get one clause instead of three, and then cannot close the proof.}}{%
\lead{Step 4 \; Resolution}
\begin{itemize}
  \item[6.] $\neg LS(x) \vee \neg Have(x,y) \vee \neg Hound(y)$
            \hfill \yr{(3,\,1) $[y/x]$}
  \item[7.] $\neg Have(John,y) \vee \neg Hound(y)$ \hfill \yr{(6,\,5a) $[John/x]$}
  \item[8.] $\neg Hound(a)$ \hfill \yr{(7,\,4a) $[a/y]$}
  \item[9.] $Cat(a)$ \hfill \yr{(8,\,4b)}
  \item[10.] $\neg Have(x,a) \vee \neg Have(x,z) \vee \neg Mouse(z)$
            \hfill \yr{(9,\,2) $[a/y]$}
  \item[11.] $\neg Have(John,z) \vee \neg Mouse(z)$ \hfill \yr{(10,\,4a) $[John/x]$}
  \item[12.] $\neg Mouse(b)$ \hfill \yr{(11,\,5b) $[b/z]$}
  \item[13.] $\square$ \hfill \yr{(12,\,5c)}
\end{itemize}
\lead{Final answer}
\begin{itemize}
  \item \mk{``If John is a light sleeper then John has no mice'' is
        proved.} \checkmark \; Eight steps.
  \item \textbf{The shape of the argument}, worth stating in words underneath: John is a
        light sleeper, so he has nothing that howls, so what he has is not a hound; it was
        a cat or a hound, so it is a cat; anyone with a cat has no mice; so he has no mice.
  \item \textbf{Two skolem constants, and they are different.} $a$ is the animal John owns
        (from premise 4); $b$ is the mouse the negated goal claims he has. Reusing one
        letter for both is a wrong proof that happens to close.
  \item \textbf{Premise 2 needs four literals.} Both the cat and the mouse are things
        \emph{John has}, so $Have$ appears twice with different second arguments.
        Writing it with one $Have$ loses the proof.
\end{itemize}}

\penalty0\vspace{2pt}

\begin{asked}{2076 Bhadra \textperiodcentered{} CT 78506 \textperiodcentered{} [7]}
Consider the following axioms: All direwolves howl at night. Anyone who has any cat will
not have any mice. Light sleepers do not have anything that howls at night. Jon has either
a cat or a direwolf. Prove that if Jon is a light sleeper, then Jon does not have any mice,
using the resolution method.
\end{asked}

{\color{sub}\footnotesize \mk{Identical}: substitute
$[Jon/John,\; Direwolf/Hound]$ and the same eight steps close the proof. Nothing else
changes.}

%=======================================================================
\T{P5. Resolution refutation \gap\textcolor{sub}{The ten one-off premise sets}}

{\color{sub}\footnotesize 11 papers, one set each except Santa which is set twice
\yr{72 Ma \m{7}, \textbf{81 Ch} \m{8}, \textbf{78 Ch} \m{3+3}, \textbf{80 Ch} \m{3+5},
\texttt{81 Ba} \m{2+6}, 79 Ash \m{7}, \textbf{\textit{77 Ch}} \m{8}, \textbf{70 Bh} \m{8},
\textbf{71 Bh} \m{8}, 73 Ma \m{10}, 80 Ash \m{3+5}}. Same five steps every time; only the
premises move.}

\creamq{1. Did Marcus hate Caesar?}{\m{7}}
{\color{sub}\footnotesize \yr{72 Ma \m{7}} \gap$\cdot$\gap the paper prints
\emph{``Find, did Marcus Caesor?''}, dropping the verb; it means \emph{hate}}

\begin{asked}{2072 Magh \textperiodcentered{} [7]}
Using resolution, solve the following statements: All pompeian are Romans. All Romans were
either loyal to Caesor or hated him. Everyone is loyal to someone. People only try to
assassinate rulers they not loyal to. Marcus tried to assassinate Caesor. Marcus was
pompeian. Find, did Marcus Caesor?
\end{asked}

\sbs{%
\lead{FOPL and CNF}
\begin{itemize}
  \item[1.] $\neg Pompeian(x) \vee Roman(x)$
  \item[2.] $\neg Roman(x) \vee loyal(x, Caesar) \vee hate(x, Caesar)$
  \item[3.] $loyal\big(x,\, f(x)\big)$ \hfill \yr{$\exists$ inside $\forall$: skolem
            \textbf{function}}
  \item[4.] $\neg tryAssassinate(x,y) \vee \neg ruler(y) \vee \neg loyal(x,y)$
  \item[5.] $tryAssassinate(Marcus, Caesar)$
  \item[6.] $Pompeian(Marcus)$ \qquad \textbf{7.} $ruler(Caesar)$
\end{itemize}
{\color{sub}\footnotesize $ruler(Caesar)$ is not printed but rule 4 is about
\emph{rulers}, so state it as an assumption. \emph{``Everyone is loyal to someone''} is
$\forall x \exists y\, loyal(x,y)$, and \mk{it skolemizes to $f(x)$, not to
a constant} $-$ each person is loyal to their own someone.}}{%
\lead{Proof \; \yr{negated goal 8. $\neg hate(Marcus, Caesar)$}}
\begin{itemize}
  \item[9.] $Roman(Marcus)$ \hfill \yr{(6,\,1)}
  \item[10.] $loyal(Marcus,Caesar) \vee hate(Marcus,Caesar)$ \hfill \yr{(9,\,2)}
  \item[11.] $loyal(Marcus, Caesar)$ \hfill \yr{(10,\,8)}
  \item[12.] $\neg ruler(Caesar) \vee \neg loyal(Marcus, Caesar)$
            \hfill \yr{(5,\,4) $[Marcus/x, Caesar/y]$}
  \item[13.] $\neg loyal(Marcus, Caesar)$ \hfill \yr{(12,\,7)}
  \item[14.] $\square$ \hfill \yr{(13,\,11)}
\end{itemize}
\begin{itemize}
  \item \mk{Marcus hated Caesar.} \checkmark
  \item \textbf{Clause 3 is never used}, and it is the only clause with a Skolem function
        in it. It is there to see whether you skolemize it correctly, not because the
        proof needs it.
\end{itemize}}

\penalty0\vspace{2pt}

\creamq{2. Scrooge is not a child}{\m{8}\m{3+3}}
{\color{sub}\footnotesize \yr{\textbf{81 Ch} \m{8}, \textbf{78 Ch} \m{3+3}}
\gap$\cdot$\gap 78 Ch prints the same six axioms and asks only for the
\textbf{CNF plus the horn-clause definition}, so the clause list below is the whole answer
there}

\begin{asked}{2081 Chaitra \textperiodcentered{} [8]}
Every child loves Santa. Everyone who loves Santa loves any reindeer. Rudolph is a
reindeer, and Rudolph has a red nose. Anything which has a red nose is weird or is a clown.
No reindeer is a clown. Scrooge does not love anything which is weird. Using Resolution
Refutation, derive the conclusion: Scrooge is not a child.
\end{asked}

\sbs{%
\lead{CNF}
\begin{itemize}
  \item[1.] $\neg child(x) \vee loves(x, Santa)$
  \item[2.] $\neg loves(x, Santa) \vee \neg reindeer(y) \vee loves(x,y)$
  \item[3.] $reindeer(Rudolph)$ \qquad \textbf{4.} $redNose(Rudolph)$
  \item[5.] $\neg redNose(x) \vee weird(x) \vee clown(x)$
  \item[6.] $\neg reindeer(x) \vee \neg clown(x)$
  \item[7.] $\neg weird(x) \vee \neg loves(Scrooge, x)$
\end{itemize}
{\color{sub}\footnotesize The goal is already a negation, so
\mk{negating it removes the $\neg$}: clause 8 is
$child(Scrooge)$, not $\neg child(Scrooge)$. Getting this backwards is the standard error
on this set.}}{%
\lead{Proof \; \yr{negated goal 8. $child(Scrooge)$}}
\begin{itemize}
  \item[9.] $loves(Scrooge, Santa)$ \hfill \yr{(8,\,1) $[Scrooge/x]$}
  \item[10.] $\neg reindeer(y) \vee loves(Scrooge, y)$ \hfill \yr{(9,\,2)}
  \item[11.] $loves(Scrooge, Rudolph)$ \hfill \yr{(10,\,3) $[Rudolph/y]$}
  \item[12.] $\neg weird(Rudolph)$ \hfill \yr{(11,\,7) $[Rudolph/x]$}
  \item[13.] $\neg redNose(Rudolph) \vee clown(Rudolph)$ \hfill \yr{(12,\,5)}
  \item[14.] $clown(Rudolph)$ \hfill \yr{(13,\,4)}
  \item[15.] $\neg reindeer(Rudolph)$ \hfill \yr{(14,\,6)}
  \item[16.] $\square$ \hfill \yr{(15,\,3)}
\end{itemize}
\begin{itemize}
  \item \mk{``Scrooge is not a child'' is proved.} \checkmark \;
        Eight steps, every premise used.
  \item The argument in words: if Scrooge were a child he would love Santa, hence love
        Rudolph; but Scrooge loves nothing weird, so Rudolph is not weird, so
        (red nose) he is a clown, and no reindeer is a clown $-$ contradiction.
\end{itemize}}

\penalty0\vspace{2pt}

\creamq{3. Did Curiosity kill the cat?}{\m{3+5}}
{\color{sub}\footnotesize \tP{1} \yr{\textbf{80 Ch} \m{3+5}} \gap$\cdot$\gap
\mk{the hardest clause set in the chapter}: premise (a) is the
Skolem-function sentence from P2, and it is needed twice}

\begin{asked}{2080 Chaitra \textperiodcentered{} [3+5]}
Convert the following sentences into FOPL and answer ``Did Curiosity kill the cat'' by
resolution refutation:
\begin{enumerate}[label=\alph*), leftmargin=1.6em]
  \item Everyone who loves all animals is loved by someone.
  \item Anyone who kills an animal is loved by no one.
  \item Jack loves all animals.
  \item Either Jack or Curiosity killed the cat.
\end{enumerate}
\end{asked}

\sbsr{0.50}{%
\lead{CNF \; \yr{(the cat is named Tuna; ``cats are animals'' is assumed)}}
\begin{itemize}
  \item[A1.] $Animal(F(x)) \vee Loves(G(x), x)$
  \item[A2.] $\neg Loves(x, F(x)) \vee Loves(G(x), x)$
  \item[B.] $\neg Animal(y) \vee \neg Kills(x,y) \vee \neg Loves(z,x)$
  \item[C.] $\neg Animal(x) \vee Loves(Jack, x)$
  \item[D.] $Kills(Jack, Tuna) \vee Kills(Curiosity, Tuna)$
  \item[E.] $Cat(Tuna)$ \qquad \textbf{F.} $\neg Cat(x) \vee Animal(x)$
\end{itemize}
{\color{sub}\footnotesize A1 and A2 are premise (a) put through all eight CNF steps in
P2 above. $F(x)$ is ``an animal $x$ does not love'' and $G(x)$ is ``someone who loves
$x$''. \mk{Both must be functions of $x$.}}}{%
\lead{Proof \; \yr{negated goal G. $\neg Kills(Curiosity, Tuna)$}}
\begin{itemize}
  \item[1.] $Kills(Jack, Tuna)$ \hfill \yr{(G,\,D)}
  \item[2.] $Animal(Tuna)$ \hfill \yr{(E,\,F)}
  \item[3.] $\neg Animal(Tuna) \vee \neg Loves(z, Jack)$
            \hfill \yr{(1,\,B) $[Jack/x, Tuna/y]$}
  \item[4.] $\neg Loves(z, Jack)$ \hfill \yr{(3,\,2)}
  \item[5.] $\neg Animal(F(Jack)) \vee Loves(G(Jack), Jack)$
            \hfill \yr{(C,\,A2) $[Jack/x,\, F(Jack)/x]$}
  \item[6.] $Loves(G(Jack), Jack)$ \hfill \yr{(5,\,A1) $[Jack/x]$, duplicate collapses}
  \item[7.] $\square$ \hfill \yr{(6,\,4) $[G(Jack)/z]$}
\end{itemize}
\begin{itemize}
  \item \mk{Curiosity killed the cat.} \checkmark
  \item \textbf{Steps 5 and 6 are the whole difficulty.} They prove
        \emph{somebody loves Jack}, using (c) to satisfy the antecedent of (a): Jack loves
        all animals, so in particular he loves $F(Jack)$, so the second clause of (a)
        fires. Only then does step 4, \emph{nobody loves Jack}, become a contradiction.
  \item \textbf{Two facts have to be assumed} because the paper omits them: the cat has a
        name, and cats are animals. State both.
\end{itemize}}

\penalty0\vspace{2pt}

\creamq{4. Sneha is happy}{\m{2+6}}
{\color{sub}\footnotesize \yr{\texttt{81 Ba} \m{2+6}} \gap$\cdot$\gap the first 2 marks
are \textbf{skolemization defined}, \S4.3 \gap$\cdot$\gap
\mk{contrast the CNF of premise (b) with the food family's premise 4}: here
the distribution into two clauses \emph{is} valid}

\begin{asked}{2081 Baishakh \textperiodcentered{} CT 710 \textperiodcentered{} [2+6]}
Define skolemization with example. Consider the following axioms: a) anyone passing his
engineering exams and winning the lottery is happy; b) anyone who studies or is lucky can
pass all his exams; c) Sneha did not study but she is lucky; d) anyone who is lucky wins
the lottery. Using resolution refutation, prove that Sneha is happy.
\end{asked}

\sbs{%
\lead{CNF}
\begin{itemize}
  \item[1.] $\neg pass(x) \vee \neg win(x) \vee happy(x)$
  \item[2.] $\big[study(x) \vee lucky(x)\big] \rightarrow pass(x)$ \;
            $= \big[\neg study(x) \wedge \neg lucky(x)\big] \vee pass(x)$, so
  \begin{itemize}
    \item[2a.] $\neg study(x) \vee pass(x)$
    \item[2b.] $\neg lucky(x) \vee pass(x)$
  \end{itemize}
  \item[3a.] $\neg study(Sneha)$ \qquad \textbf{3b.} $lucky(Sneha)$
  \item[4.] $\neg lucky(x) \vee win(x)$
\end{itemize}
{\color{sub}\footnotesize \mk{Why this splits and the food premise does
not}: here the negation lands on a \emph{disjunction},
$\neg(A \vee B) = \neg A \wedge \neg B$, so distributing is forced. In the food family it
lands on a \emph{conjunction}, $\neg(A \wedge \neg B) = \neg A \vee B$, and there is
nothing to split. One De Morgan apart, opposite answers.}}{%
\lead{Proof \; \yr{negated goal 5. $\neg happy(Sneha)$}}
\begin{itemize}
  \item[6.] $\neg pass(Sneha) \vee \neg win(Sneha)$ \hfill \yr{(5,\,1) $[Sneha/x]$}
  \item[7.] $\neg lucky(Sneha) \vee \neg win(Sneha)$ \hfill \yr{(6,\,2b)}
  \item[8.] $\neg win(Sneha)$ \hfill \yr{(7,\,3b)}
  \item[9.] $\neg lucky(Sneha)$ \hfill \yr{(8,\,4) $[Sneha/x]$}
  \item[10.] $\square$ \hfill \yr{(9,\,3b)}
\end{itemize}
\begin{itemize}
  \item \mk{``Sneha is happy'' is proved.} \checkmark \; Five steps.
  \item \textbf{Clauses 2a and 3a are never used.} \emph{``Sneha did not study''} is there
        only to tempt you into deriving nothing from it; the lucky half carries the whole
        proof.
  \item \emph{``can pass all his exams''} is written $pass(x)$ with the exam left implicit.
        Adding a second argument $pass(x, exam)$ changes nothing except the typing.
\end{itemize}}

\penalty0\vspace{2pt}

\creamq{5. Steve likes BK301}{\m{7}}
{\color{sub}\footnotesize \yr{79 Ash \m{7}} \gap$\cdot$\gap
\mk{the printed premise runs the wrong way} and has to be read as its
converse before anything can be proved}

\begin{asked}{2079 Ashwin \textperiodcentered{} [7]}
Assume the following facts: a) Steve only likes easy courses; b) Science courses are hard;
c) All the courses in the basket weaving department are easy; d) BK301 is a basket weaving
course. Prove that Steve likes BK301 using Resolution Refutation.
\end{asked}

\sbs{%
\lead{\mk{Read premise (a) carefully first}}
\begin{itemize}
  \item Literally, \emph{``Steve only likes easy courses''} is
        $\forall x\; likes(Steve,x) \rightarrow easy(x)$: liking implies easy.
  \item From that you can prove \emph{a course Steve likes is easy}. You cannot prove he
        likes anything, because nothing in the set ever concludes $likes$.
  \item The question is only answerable under the converse reading,
        \emph{``Steve likes every easy course''}:
        \[ \forall x\; easy(x) \rightarrow likes(Steve, x) \]
  \item \mk{Write that line, then prove it.} Stating the reading you
        adopted is a mark; silently reversing an implication is not.
\end{itemize}}{%
\lead{CNF and proof}
\begin{itemize}
  \item[1.] $\neg easy(x) \vee likes(Steve, x)$ \hfill \yr{(a), converse reading}
  \item[2.] $\neg science(x) \vee \neg easy(x)$ \hfill \yr{(b)}
  \item[3.] $\neg basketWeaving(x) \vee easy(x)$ \hfill \yr{(c)}
  \item[4.] $basketWeaving(BK301)$ \hfill \yr{(d)}
  \item[5.] $\neg likes(Steve, BK301)$ \hfill \yr{negated goal}
  \item[6.] $\neg easy(BK301)$ \hfill \yr{(5,\,1) $[BK301/x]$}
  \item[7.] $\neg basketWeaving(BK301)$ \hfill \yr{(6,\,3)}
  \item[8.] $\square$ \hfill \yr{(7,\,4)}
\end{itemize}
\begin{itemize}
  \item \mk{``Steve likes BK301'' is proved} under that reading.
        \checkmark \; Three steps.
  \item \textbf{Premise (b) is unused}, and it is the only one mentioning science.
        \emph{Hard} is written $\neg easy$; inventing a separate $hard$ predicate needs a
        fifth premise saying the two are opposites.
\end{itemize}}

\penalty0\vspace{2pt}

\creamq{6. You will enjoy}{\m{8}}
{\color{sub}\footnotesize \yr{\textbf{\textit{77 Ch}} \m{8}} \gap$\cdot$\gap
purely \textbf{propositional}, no quantifiers, and \mk{four of the seven
premises are noise}}

\begin{asked}{2077 Chaitra \textperiodcentered{} CT 78506 \textperiodcentered{} [8]}
If it is a sunny and warm day you will enjoy. If it is a warm and pleasant day you will do
strawberry picking. If it is raining, no strawberry picking. If it is raining you will get
wet. It is a warm day. It is raining. It is sunny. Prove that ``You will enjoy'' using
resolution by refutation.
\end{asked}

\sbs{%
\lead{CNF}
\begin{itemize}
  \item[1.] $\neg sunny \vee \neg warm \vee enjoy$
  \item[2.] $\neg warm \vee \neg pleasant \vee picking$
  \item[3.] $\neg raining \vee \neg picking$
  \item[4.] $\neg raining \vee wet$
  \item[5.] $warm$ \qquad \textbf{6.} $raining$ \qquad \textbf{7.} $sunny$
\end{itemize}
{\color{sub}\footnotesize Each implication becomes one clause by
$\alpha \rightarrow \beta = \neg\alpha \vee \beta$ plus De Morgan on the conjunctive
antecedent. No quantifiers means no unification: every resolution step here is on a
\textbf{ground} literal.}}{%
\lead{Proof \; \yr{negated goal 8. $\neg enjoy$}}
\begin{itemize}
  \item[9.] $\neg sunny \vee \neg warm$ \hfill \yr{(8,\,1)}
  \item[10.] $\neg warm$ \hfill \yr{(9,\,7)}
  \item[11.] $\square$ \hfill \yr{(10,\,5)}
\end{itemize}
\begin{itemize}
  \item \mk{``You will enjoy'' is proved} in three steps. \checkmark
  \item \textbf{Clauses 2, 3, 4 and 6 are never touched.} The strawberry picking and the
        rain are a complete decoy: they are consistent with the rest and irrelevant to the
        goal.
  \item \mk{The lesson worth writing under the proof}: resolve
        \emph{from the negated goal}. Starting from the facts, this set will happily
        generate $wet$, $\neg picking$ and several other true-but-useless literals before
        it stumbles on the answer.
\end{itemize}}

\penalty0\vspace{2pt}

\creamq{7. Ram is naughty}{\m{8}}
{\color{sub}\footnotesize \yr{\textbf{70 Bh} \m{8}} \gap$\cdot$\gap also set as an
exercise by the textbook, whose own solution omits the argument order}

\begin{asked}{2070 Bhadra \textperiodcentered{} [8]}
All oversmart persons are stupid. Children of oversmart persons are naughty. Ram is
children of Hari. Hari is oversmart. Show that Ram is naughty using FOPL based resolution.
\end{asked}

\sbs{%
\lead{CNF}
\begin{itemize}
  \item[1.] $\neg oversmart(x) \vee stupid(x)$
  \item[2.] $\neg child(x,y) \vee \neg oversmart(y) \vee naughty(x)$
  \item[3.] $child(Ram, Hari)$ \qquad \textbf{4.} $oversmart(Hari)$
\end{itemize}
\lead{Proof \; \yr{negated goal 5. $\neg naughty(Ram)$}}
\begin{itemize}
  \item[6.] $\neg child(Ram, y) \vee \neg oversmart(y)$ \hfill \yr{(5,\,2) $[Ram/x]$}
  \item[7.] $\neg oversmart(Hari)$ \hfill \yr{(6,\,3) $[Hari/y]$}
  \item[8.] $\square$ \hfill \yr{(7,\,4)}
\end{itemize}}{%
\lead{Notes}
\begin{itemize}
  \item \mk{``Ram is naughty'' is proved} in three steps. \checkmark
  \item \textbf{Premise 1 is unused.} \emph{Stupid} appears nowhere else, so the whole
        first sentence is a decoy $-$ the same trick as the mammal rule in Family B.
  \item \textbf{$child(x,y)$ must mean ``$x$ is a child of $y$''} and be used that way in
        both clause 2 and clause 3. The English \emph{``Ram is children of Hari''} fixes
        the order; write the reading down before you write the clause, because reversing
        it silently breaks the unification in step 6.
  \item Predicate names are free, but $naughty$ must be a \emph{property of the child}, not
        a relation: $naughty(x)$, one argument.
\end{itemize}}

\penalty0\vspace{2pt}

\creamq{8. The book supports the cup}{\m{8}}
{\color{sub}\footnotesize \yr{\textbf{71 Bh} \m{8}} \gap$\cdot$\gap the only set written
in \textbf{relational} rather than sentence form, so the argument order is the whole
question}

\begin{asked}{2071 Bhadra \textperiodcentered{} [8]}
If X is on the top of Y, Y supports X. If X is above Y and they are touching each other, X
is on the top of Y. A cup is above a book. A cup is touching a book. Show that
supports(book, cup) is true.
\end{asked}

\sbs{%
\lead{CNF}
\begin{itemize}
  \item[1.] $\neg onTop(x,y) \vee supports(y,x)$
  \item[2.] $\neg above(x,y) \vee \neg touching(x,y) \vee onTop(x,y)$
  \item[3.] $above(Cup, Book)$ \qquad \textbf{4.} $touching(Cup, Book)$
\end{itemize}
\lead{Proof \; \yr{negated goal 5. $\neg supports(Book, Cup)$}}
\begin{itemize}
  \item[6.] $\neg onTop(Cup, Book)$ \hfill \yr{(5,\,1) $[Cup/x, Book/y]$}
  \item[7.] $\neg above(Cup,Book) \vee \neg touching(Cup,Book)$ \hfill \yr{(6,\,2)}
  \item[8.] $\neg touching(Cup, Book)$ \hfill \yr{(7,\,3)}
  \item[9.] $\square$ \hfill \yr{(8,\,4)}
\end{itemize}}{%
\lead{Notes}
\begin{itemize}
  \item \mk{$supports(Book, Cup)$ is proved} in four steps. \checkmark
  \item \mk{Clause 1 reverses its arguments and that is the point.}
        $onTop(x,y)$ has the supporter second; $supports(y,x)$ has it first. The goal
        $supports(Book,Cup)$ therefore unifies with $[Cup/x, Book/y]$, giving
        $onTop(Cup,Book)$ $-$ cup on book, which is what the facts describe.
  \item \textbf{Every premise is used}, and there is no quantifier subtlety: the two rules
        are plain universals, so drop the $\forall$ and go.
  \item \emph{``they are touching each other''} is symmetric in English but
        $touching(x,y)$ is not symmetric in the clause set. It does not matter here
        because the facts are given in the same order the rule wants; if they were not,
        a symmetry axiom $\neg touching(x,y) \vee touching(y,x)$ would be needed.
\end{itemize}}

\penalty0\vspace{2pt}

\creamq{9. If John does not study, Mary does not love John}{\m{10}}
{\color{sub}\footnotesize \tP{1} \yr{73 Ma \m{10}} \gap$\cdot$\gap the goal is an
implication again, so \mk{negating it gives two clauses}, and every premise
is a contrapositive}

\begin{asked}{2073 Magh \textperiodcentered{} [10]}
Consider the following axioms: i) anyone whom Mary loves is a football star; ii) any
student who does not pass does not play; iii) John is a student; iv) any student who does
not study does not pass; v) anyone who does not play is not a football star. Prove that
``If John does not study, then Mary does not love John'' using resolution by refutation.
\end{asked}

\sbsr{0.48}{%
\lead{CNF}
\begin{itemize}
  \item[1.] $\neg loves(Mary,x) \vee star(x)$
  \item[2.] $\neg student(x) \vee pass(x) \vee \neg play(x)$
  \item[3.] $student(John)$
  \item[4.] $\neg student(x) \vee study(x) \vee \neg pass(x)$
  \item[5.] $play(x) \vee \neg star(x)$
\end{itemize}
\lead{\mk{Negate the goal: two clauses}}
\begin{itemize}
  \item $\neg\big[\neg study(John) \rightarrow \neg loves(Mary, John)\big]
        = \neg study(John) \wedge loves(Mary, John)$
  \item[6a.] $\neg study(John)$ \qquad \textbf{6b.} $loves(Mary, John)$
\end{itemize}
{\color{sub}\footnotesize Premises (ii), (iv) and (v) are all of the form
\emph{not-$P$ implies not-$Q$}. Convert each one mechanically:
$[\,S \wedge \neg P\,] \rightarrow \neg Q$ becomes
$\neg S \vee P \vee \neg Q$. Trying to simplify them into positive form first is where
this question is usually lost.}}{%
\lead{Proof}
\begin{itemize}
  \item[7.] $star(John)$ \hfill \yr{(6b,\,1) $[John/x]$}
  \item[8.] $play(John)$ \hfill \yr{(7,\,5) $[John/x]$}
  \item[9.] $\neg student(John) \vee pass(John)$ \hfill \yr{(8,\,2) $[John/x]$}
  \item[10.] $pass(John)$ \hfill \yr{(9,\,3)}
  \item[11.] $\neg student(John) \vee study(John)$ \hfill \yr{(10,\,4) $[John/x]$}
  \item[12.] $study(John)$ \hfill \yr{(11,\,3)}
  \item[13.] $\square$ \hfill \yr{(12,\,6a)}
\end{itemize}
\begin{itemize}
  \item \mk{The implication is proved.} \checkmark \; Seven steps,
        every premise used, and clause 3 used \textbf{twice}.
  \item The chain in words: suppose Mary loves John and he does not study. Loved by Mary
        $\rightarrow$ football star $\rightarrow$ plays $\rightarrow$ (being a student)
        passes $\rightarrow$ studies. Contradiction with \emph{does not study}.
  \item \mk{A clause may be reused as often as needed.} Students often
        cross a clause off after one use; nothing in resolution says that.
\end{itemize}}

\penalty0\vspace{2pt}

\creamq{10. Can anyone be found with an exciting life?}{\m{3+5}}
{\color{sub}\footnotesize \yr{80 Ash \m{3+5}} \gap$\cdot$\gap asks \emph{who}, not
\emph{prove}, so this is the chapter's one genuine
\mk{answer-extraction} question}

\begin{asked}{2080 Ashwin \textperiodcentered{} [3+5]}
List the rules of inference. Convert the following into FOPL and prove using Resolution
Refutation System: ``All people who are not poor and are smart are happy. Those people who
can read are not stupid. John can read and is wealthy. Happy people have exciting lives.
Can anyone be found with an exciting life?''
\end{asked}

\sbs{%
\lead{CNF \; \yr{with \emph{wealthy} $= \neg poor$ and \emph{smart} $= \neg stupid$}}
\begin{itemize}
  \item[1.] $poor(x) \vee \neg smart(x) \vee happy(x)$
  \item[2.] $\neg read(x) \vee \neg stupid(x)$
  \item[3a.] $read(John)$ \qquad \textbf{3b.} $\neg poor(John)$
  \item[4.] $\neg happy(x) \vee exciting(x)$
  \item[5.] $\neg smart(x) \vee \neg stupid(x)$ is \textbf{not} needed; instead state the
        bridging fact \; $\neg stupid(x) \rightarrow smart(x)$, i.e.\
        \textbf{5.} $stupid(x) \vee smart(x)$
\end{itemize}
{\color{sub}\footnotesize \mk{Two vocabulary bridges must be stated}:
\emph{wealthy} is $\neg poor$, and \emph{not stupid} is \emph{smart}. Without them the
clause set has five predicates that never meet and no proof exists. Say so in one line
before starting $-$ the examiner's own wording forces it.}}{%
\lead{Proof \; \yr{negated goal 6. $\neg exciting(x)$, for all $x$}}
\begin{itemize}
  \item[7.] $\neg happy(x)$ \hfill \yr{(6,\,4)}
  \item[8.] $poor(x) \vee \neg smart(x)$ \hfill \yr{(7,\,1)}
  \item[9.] $\neg smart(John)$ \hfill \yr{(8,\,3b) $[John/x]$}
  \item[10.] $stupid(John)$ \hfill \yr{(9,\,5) $[John/x]$}
  \item[11.] $\neg read(John)$ \hfill \yr{(10,\,2) $[John/x]$}
  \item[12.] $\square$ \hfill \yr{(11,\,3a)}
\end{itemize}
\begin{itemize}
  \item \mk{Yes: John has an exciting life.} \checkmark
  \item \textbf{How the name is recovered}, which is what \emph{``can anyone be found''}
        asks: the goal $\exists x\, exciting(x)$ negates to
        $\forall x\, \neg exciting(x)$, one clause with a free variable. The
        substitutions used along the way bound that variable to \textbf{John}, and that
        binding \emph{is} the answer.
  \item The formal device is an \textbf{answer literal}: carry
        $exciting(x) \vee Ans(x)$ instead of the bare negation, and read the answer off
        the clause that survives. Worth naming even if you do not use it.
\end{itemize}}

%=======================================================================
\T{P7. Bayes' theorem: the twelve numericals}

{\color{sub}\footnotesize \mk{15 papers set a Bayes question}; twelve give
numbers and are worked below, three ask for a worked example of your own and are answered
by any one of them. \textbf{Every answer here was recomputed independently in
\code{verify.py}}, not copied from a source. The recipe never changes:}

\lead{The recipe}
\begin{enumerate}
  \item \textbf{Name the events} and write down which probability each given number is.
        \mk{Half the lost marks happen here}, in the first thirty seconds.
  \item Write Bayes' theorem with the names, not with $A$ and $B$.
  \item If $P(\text{evidence})$ is not given, build it with the law of total probability:
        $P(E) = P(E|H)P(H) + P(E|\neg H)P(\neg H)$.
  \item Substitute, and quote the answer as a percentage as well as a decimal.
\end{enumerate}

\penalty0\vspace{2pt}

\creamq{B1--B2. Meningitis and the stiff neck}{\m{3+5}\m{2+6}}
{\color{sub}\footnotesize \yr{71 Ma \m{3+5}, \texttt{81 Ba} \m{2+6}} \gap$\cdot$\gap
same scenario, \mk{different numbers in each paper} $-$ read them off the
paper, not from memory}

\begin{asked}{2071 Magh \textperiodcentered{} [3+5]}
What is neural network? Describe its types. A doctor knows that the disease meningitis
causes the patient to have a stiff neck 50\% of the time. The doctor also knows that the
probability that a patient has meningitis is 1/50,000, and the probability that any patient
has a stiff neck is 1/20. Find the probability that a patient with a stiff neck has
meningitis?
\end{asked}

\sbs{%
\lead{B1 \; 2071 Magh}
\begin{itemize}
  \item \textbf{Given} \; $P(s|m) = 0.5$, \; $P(m) = 1/50000 = 2 \times 10^{-5}$, \;
        $P(s) = 1/20 = 0.05$.
  \item \textbf{Wanted} \; $P(m|s)$.
  \item \[ P(m|s) = \frac{P(s|m)\,P(m)}{P(s)}
        = \frac{0.5 \times 2 \times 10^{-5}}{0.05}
        = \frac{10^{-5}}{0.05} \]
  \item \[ \boxed{P(m|s) = 0.0002 = 0.02\,\%} \]
  \item \mk{$P(s)$ is given here}, so no total-probability step is
        needed. That is the only easy version of this question in the archive.
\end{itemize}}{%
\lead{B2 \; 2081 Baishakh \yr{(CT 710)}}
\begin{itemize}
  \item \textbf{Given} \; $P(s|m) = 0.8$, \; $P(m) = 1/50000$, \; $P(s) = 0.01$.
  \item \[ P(m|s) = \frac{0.8 \times 2 \times 10^{-5}}{0.01}
        = \frac{1.6 \times 10^{-5}}{0.01} \]
  \item \[ \boxed{P(m|s) = 0.0016 = 0.16\,\%} \]
  \item The first 2 marks are \textbf{prior Vs posterior} defined: $P(m)$ is the prior,
        $P(m|s)$ the posterior, and the stiff neck raised the belief from
        \mk{0.002\,\% to 0.16\,\%} $-$ eighty times higher, and still
        under one in six hundred.
  \item \mk{That last sentence is the exam answer to ``what does this
        tell you''}: a strong likelihood cannot outweigh a tiny prior.
\end{itemize}}

\penalty0\vspace{2pt}

\creamq{B3. Measles or flu?}{\m{8}}
{\color{sub}\footnotesize \yr{73 Ma \m{8}} \gap$\cdot$\gap the priors are given directly
as a 90/10 split, so this is the cleanest total-probability example in the archive}

\begin{asked}{2073 Magh \textperiodcentered{} [8]}
A doctor is called to see a sick child. The doctor has prior information that 90\% of sick
children in that neighborhood have the flu, while the other 10\% are sick with measles. Let
F stand for an event of a child being sick with Flu and M stand for an event of a child
being sick with measles. Assume for simplicity that there no other maladies in that
neighborhood. A well-known (and common) symptom of measles is a rash and has probability of
0.95. However, very occasionally, children with flu also develop rash and has probability
of 0.08. Upon examining the child, the doctor finds a rash. What is the probability that the
child has measles?
\end{asked}

\sbs{%
\lead{Solution}
\begin{itemize}
  \item \textbf{Given} \; $P(F) = 0.9$, $P(M) = 0.1$, $P(R|M) = 0.95$, $P(R|F) = 0.08$.
  \item \textbf{Total probability of a rash}:
        \[ P(R) = P(R|M)P(M) + P(R|F)P(F) \]
        \[ = 0.95 \times 0.1 + 0.08 \times 0.9 = 0.095 + 0.072 = 0.167 \]
  \item \[ P(M|R) = \frac{P(R|M)P(M)}{P(R)} = \frac{0.095}{0.167} \]
  \item \[ \boxed{P(M|R) = 0.5689 \approx 56.9\,\%} \]
\end{itemize}}{%
\lead{Reading the answer}
\begin{itemize}
  \item Measles goes from a \textbf{10\,\% prior to a 57\,\% posterior}: the rash is
        strong evidence, but flu is nine times commoner, so it stays a coin flip.
  \item \mk{Check the complement}:
        $P(F|R) = 0.072/0.167 = 0.4311$, and $0.5689 + 0.4311 = 1$. \checkmark \;
        Two extra lines, and they catch an arithmetic slip in the denominator every time.
  \item \textbf{The trap}: reading 0.95 as $P(M|R)$ and answering ``95\,\%''. The paper
        gives the \emph{likelihood} $P(R|M)$; the whole point of the question is that the
        two are not the same number.
\end{itemize}}

\penalty0\vspace{2pt}

\creamq{B4. The cigar smoker}{\m{2+6}\m{8}}
{\color{sub}\footnotesize \yr{\textbf{\texttt{80 Bh}} \m{2+6}, \textbf{\textit{73 Bh}}
\m{8}} \gap$\cdot$\gap two papers, identical data \gap$\cdot$\gap part (i) is the
\textbf{prior}, and it is one mark for writing down a number already in the question}

\begin{asked}{2080 Bhadra \textperiodcentered{} CT 710 \textperiodcentered{} [2+6] \; also 2073 Bhadra CT 78506 [8]}
What do you understand by Bayesian Network? In a County, 51\% of the adults are males and
the other 49\% are females. One adult is randomly selected for a survey involving credit
card usage. i) Find the prior probability that the selected person is a male. ii) It is
later learned that the selected survey subject was smoking a cigar. Also, 9.5\% of males
smoke cigars, whereas 1.7\% of females smoke cigars (based on data from the Substance Abuse
and Mental Health Services Administration). Use this additional information to find the
probability that the selected subject is a male.
\end{asked}

\sbs{%
\lead{Solution}
\begin{itemize}
  \item \textbf{(i) Prior} \; $P(M) = 0.51$. \; \mk{That is the whole
        answer to part (i)}; no calculation is intended.
  \item \textbf{(ii) Given} \; $P(F) = 0.49$, $P(C|M) = 0.095$, $P(C|F) = 0.017$.
  \item \[ P(C) = 0.095 \times 0.51 + 0.017 \times 0.49 \]
        \[ = 0.04845 + 0.00833 = 0.05678 \]
  \item \[ P(M|C) = \frac{0.04845}{0.05678} \]
  \item \[ \boxed{P(M|C) = 0.8533 \approx 85.3\,\%} \]
\end{itemize}}{%
\lead{Notes}
\begin{itemize}
  \item The prior 51\,\% rises to \textbf{85\,\%} because men are five and a half times as
        likely to smoke a cigar. \mk{Here the likelihood ratio does
        dominate}, which is the opposite lesson to B1 and B2 $-$ and the
        difference is entirely that this prior is near a half rather than near zero.
  \item \textbf{Keep five decimal places in the denominator.} Rounding $0.05678$ to
        $0.057$ moves the answer to 85.0\,\%, and the marks for the final figure are
        awarded on the exact value.
  \item \yr{\textbf{\textit{73 Bh}}} prints the same scenario for a single 8-mark part.
        Same two answers.
\end{itemize}}

\penalty0\vspace{2pt}

\creamq{B5. The 99\,\% accurate test for a rare disease}{\m{8}}
{\color{sub}\footnotesize \tP{1} \yr{\textbf{\textit{72 Ash}} \m{8}} \gap$\cdot$\gap
\mk{the base-rate question}, and the one every paper's theory part is
really about}

\begin{asked}{2072 Ashwin \textperiodcentered{} CT 78506 \textperiodcentered{} [8]}
After your yearly checkup, the doctor has bad news and good news. The bad news is that you
tested positive for a serious disease, and that the test is 99\% accurate (i.e.\ the
probability of testing positive given that you have the disease is 0.99, as is the
probability of testing negative given that you don't have the disease). The good news is
that this is a rare disease, striking only one in 10,000 people. Why is it good news that
the disease is rare? What are the chances that you actually have the disease?
\end{asked}

\sbs{%
\lead{Solution}
\begin{itemize}
  \item \textbf{Given} \; $P(D) = 10^{-4}$, $P(+|D) = 0.99$, $P(-|\neg D) = 0.99$, so
        \mk{$P(+|\neg D) = 0.01$}.
  \item \[ P(+) = 0.99 \times 10^{-4} + 0.01 \times 0.9999 \]
        \[ = 0.000099 + 0.009999 = 0.010098 \]
  \item \[ P(D|+) = \frac{0.000099}{0.010098} \]
  \item \[ \boxed{P(D|+) = 0.0098 \approx 0.98\,\%} \]
  \item \textbf{Fewer than 1 in 100} people with a positive result actually have the
        disease.
\end{itemize}}{%
\lead{\mk{``Why is it good news that the disease is rare''}}
\begin{itemize}
  \item Because the \textbf{false positives outnumber the true ones}, and by a lot.
        In 1,000,000 people: \textbf{100} have the disease and 99 of them test positive;
        999,900 do not, and \textbf{9,999} of them still test positive.
        \mk{9,999 against 99 is roughly 100 to 1.}
  \item So a rare disease means a positive result is far more likely to be an error than
        a diagnosis. Rarity protects you.
  \item \textbf{Write the population count.} It is worth more marks than the algebra and
        it is much harder to get wrong.
  \item \textbf{The general lesson}: the posterior depends on the prior as much as on the
        test. A 99\,\% test is not a 99\,\% answer.
\end{itemize}}

\penalty0\vspace{2pt}

\creamq{B6. The genetic defect}{\m{2+6}}
{\color{sub}\footnotesize \yr{81 Ash \m{2+6}} \gap$\cdot$\gap part (b) asks for the
\textbf{probability that they do \emph{not} have it}, which is the complement $-$
\mk{read the question twice}}

\begin{asked}{2081 Ashwin \textperiodcentered{} [2+6]}
What are the approaches for measuring uncertainty? 1\% of people have a certain genetic
defect. 90\% of tests for the gene detect the defect (true positives). 10\% of the tests
are false positives. a) What might be the probability of getting the genetic defect result?
b) If a person gets a positive test result, what are the odds they actually don't have the
genetic defect?
\end{asked}

\sbs{%
\lead{Solution}
\begin{itemize}
  \item \textbf{Given} \; $P(D) = 0.01$, $P(+|D) = 0.90$, $P(+|\neg D) = 0.10$.
  \item \textbf{(a)} The probability of a positive result:
        \[ P(+) = 0.90 \times 0.01 + 0.10 \times 0.99 \]
        \[ = 0.009 + 0.099 = \boxed{0.108 = 10.8\,\%} \]
  \item \textbf{(b)} \[ P(\neg D \,|\, +) = \frac{0.099}{0.108} \]
        \[ = \boxed{0.9167 \approx 91.7\,\%} \]
  \item And so $P(D|+) = 0.009/0.108 = 0.0833 = 8.3\,\%$. \;
        $91.7 + 8.3 = 100$. \checkmark
\end{itemize}}{%
\lead{Notes}
\begin{itemize}
  \item \textbf{As true odds}, which is what part (b) literally asks for:
        $0.099 : 0.009$, i.e.\ \mk{11 to 1 against actually having the
        defect}. Give both the probability and the odds and neither reading
        can cost a mark.
  \item \textbf{Part (a) is ambiguous} as printed $-$ \emph{``the probability of getting
        the genetic defect result''} could mean $P(D)$, which is stated in the question
        and would make the part worthless. It means $P(+)$. Say which reading you took.
  \item Same shape as B5 with a much weaker test (10\,\% false positives instead of 1\,\%)
        and a much commoner condition (1\,\% instead of 0.01\,\%). The posterior lands at
        8\,\% rather than 1\,\%.
\end{itemize}}

\penalty0\vspace{2pt}

\creamq{B7. Which factory made the defective bulb?}{\m{3+5}}
{\color{sub}\footnotesize \yr{80 Ash \m{3+5}} \gap$\cdot$\gap
\mk{three-way partition}, so the denominator has three terms}

\begin{asked}{2080 Ashwin \textperiodcentered{} [3+5]}
How is reasoning done under uncertainty? Three factories produce light bulbs to supply the
market. Factory A produces 20\%, 50\% of the bulbs are produced in factory B and 30\% in
factory C. 2\% of the bulbs produced in factory A, 5\% of the bulbs produced in factory B
and 3\% of the bulbs produced in factory C are defective. A bulb is selected at random in
the market and found to be defective. What is the probability that this bulb was produced by
factory B?
\end{asked}

\sbs{%
\lead{Solution}
\begin{itemize}
  \item \textbf{Given} \; $P(A) = 0.2$, $P(B) = 0.5$, $P(C) = 0.3$; \;
        $P(D|A) = 0.02$, $P(D|B) = 0.05$, $P(D|C) = 0.03$.
  \item \textbf{Total probability of a defective bulb}, one term per factory:
        \[ P(D) = 0.2(0.02) + 0.5(0.05) + 0.3(0.03) \]
        \[ = 0.004 + 0.025 + 0.009 = 0.038 \]
  \item \[ P(B|D) = \frac{P(D|B)P(B)}{P(D)} = \frac{0.025}{0.038} \]
  \item \[ \boxed{P(B|D) = 0.6579 \approx 65.8\,\%} \]
\end{itemize}}{%
\lead{Notes}
\begin{itemize}
  \item \textbf{The partition must be exhaustive}: $0.2 + 0.5 + 0.3 = 1$. \checkmark \;
        Check it before you start; a partition that does not sum to 1 means a misread
        question.
  \item Factory B's share rises from \textbf{50\,\% to 66\,\%} because it is both the
        biggest producer \emph{and} the worst. The other two posteriors are
        $P(A|D) = 0.004/0.038 = 10.5\,\%$ and $P(C|D) = 0.009/0.038 = 23.7\,\%$, and the
        three sum to 100\,\%. \checkmark
  \item \mk{This is the template for every ``which source did it come
        from'' question}, and there are three of them in the archive (bulbs,
        factories, and the tall-student one below).
\end{itemize}}

\penalty0\vspace{2pt}

\creamq{B8--B9. The tall student, and the mammogram}{\m{3+5}\m{2+6}}
{\color{sub}\footnotesize \yr{\textbf{78 Ch} \m{3+5}, \texttt{82 Ba} \m{2+6}}}

\begin{asked}{2078 Chaitra \textperiodcentered{} [3+5]}
Explain how statistical reasoning aids in inference and reasoning in light of Bayes'
theorem. At a certain University, 5\% of men are over 6 feet tall and 2\% of women are 6
feet tall. 60\% of students are female. If a student is selected at random from among all
those over six feet tall, what is the probability that the selected student is a woman?
\end{asked}

\begin{asked}{2082 Baishakh \textperiodcentered{} CT 710 \textperiodcentered{} [2+6]}
What are casual networks? Given the following statistics, what is the probability that a
woman has cancer if she has a positive mammogram result? i) One percent of women over 50
have breast cancer. ii) Ninety percent of women who have breast cancer test positive on
mammograms. iii) Eight percent of women will have false positives.
\end{asked}

\sbs{%
\lead{B8 \; the tall student}
\begin{itemize}
  \item \textbf{Given} \; $P(W) = 0.6$, $P(M) = 0.4$, $P(T|W) = 0.02$, $P(T|M) = 0.05$.
  \item \[ P(T) = 0.02(0.6) + 0.05(0.4) = 0.012 + 0.020 = 0.032 \]
  \item \[ P(W|T) = \frac{0.012}{0.032} = \boxed{0.375 = 37.5\,\%} \]
  \item Women are the \textbf{majority} of students and the \textbf{minority} of tall
        ones: 60\,\% prior, 37.5\,\% posterior. \mk{The evidence reverses
        the balance}, which is the sentence the 3-mark theory part wants.
\end{itemize}}{%
\lead{B9 \; the mammogram}
\begin{itemize}
  \item \textbf{Given} \; $P(C) = 0.01$, $P(+|C) = 0.90$, $P(+|\neg C) = 0.08$.
  \item \[ P(+) = 0.90(0.01) + 0.08(0.99) \]
        \[ = 0.009 + 0.0792 = 0.0882 \]
  \item \[ P(C|+) = \frac{0.009}{0.0882} = \boxed{0.1020 \approx 10.2\,\%} \]
  \item \mk{Nine positives in ten are false alarms}, even though the
        test catches 90\,\% of real cancers. Same base-rate lesson as B5, in the setting
        it was originally made famous in.
  \item \textbf{``Eight percent of women will have false positives''} means
        $P(+|\neg C) = 0.08$, not $P(\neg C \wedge +) = 0.08$. The looser reading gives
        $0.009/0.089 = 10.1\,\%$, close enough to hide the error $-$
        \mk{state your reading}.
\end{itemize}}

\penalty0\vspace{2pt}

%=======================================================================
\T{P8. Bayesian networks: reading a joint probability off the graph}

\creamq{B10. The wet grass network}{\m{2+6}}
{\color{sub}\footnotesize \tP{1} \yr{78 Ba \m{2+6}} \gap$\cdot$\gap the paper prints the
network and its four CPTs \gap$\cdot$\gap \mk{one multiplication, four
factors}}

\begin{asked}{2078 Baishakh \textperiodcentered{} [2+6]}
What do you understand by Bayesian Network? For the given Bayesian network below, find the
probability of grass being wet when the weather is cloudy, there is some sprinkler but no
rain, i.e.\ C = True, S = True, R = False and W = True.
\end{asked}

\figC{c4_78ba_bbn.png}{0.72\textwidth}

\sbs{%
\lead{Solution}
\begin{itemize}
  \item \textbf{The product form}, one factor per node, each conditioned on its parents:
        \[ P(C, S, \neg R, W) = P(C)\,P(S|C)\,P(\neg R|C)\,P(W|S, \neg R) \]
  \item \textbf{Read each factor off its own table}:
  \begin{itemize}
    \item $P(C) = 0.6$ \hfill \yr{root node, prior}
    \item $P(S|C{=}T) = 0.10$ \hfill \yr{sprinkler table, row T}
    \item $P(R|C{=}T) = 0.80$, so
          \mk{$P(\neg R|C{=}T) = 1 - 0.80 = 0.20$}
    \item $P(W|S{=}T, R{=}F) = 0.90$ \hfill \yr{wet-grass table, row T F}
  \end{itemize}
  \item \[ = 0.6 \times 0.10 \times 0.20 \times 0.90 \]
  \item \[ \boxed{P = 0.0108 = 1.08\,\%} \]
\end{itemize}}{%
\lead{\mk{The three ways this goes wrong}}
\begin{itemize}
  \item \textbf{Forgetting to complement $R$.} The query says \emph{no rain}, and the table
        gives $P(R|C)$. Take $1 - 0.80$. Using 0.80 gives 0.0432, four times too big, and
        it is the commonest wrong answer.
  \item \textbf{Picking the wrong CPT row.} The wet-grass table is indexed by
        \textbf{(S, R)}, and the query is $S = T$, $R = F$: row \textbf{T F}, value 0.90.
        Row F T is also 0.90, so a misread here happens to give the right number in this
        one query and the wrong one in every other.
  \item \textbf{Conditioning on the wrong parent.} Sprinkler and Rain both depend on
        \emph{Cloudy} only; \emph{WetGrass} depends on both of them and \emph{not} on
        Cloudy. That is the conditional independence the network encodes, and it is what
        makes four numbers enough.
  \item \textbf{The 2-mark part}: a Bayesian network is a DAG of random variables with a
        CPT at each node, and the joint distribution is the product above. Six variables
        of this kind would need 31 free numbers as a full joint table; this one needs 9.
\end{itemize}}

\penalty0\vspace{2pt}

\creamq{B11. Smoker $\rightarrow$ cancer $\rightarrow$ test}{\m{7}}
{\color{sub}\footnotesize \tP{1} \yr{\textbf{\texttt{82 Bh}} \m{7}} \gap$\cdot$\gap
asks for an \emph{explanation}, but the marks are safest if the explanation carries
numbers \gap$\cdot$\gap a \mk{chain} network, so every inference is a
two-step}

\begin{asked}{2082 Bhadra \textperiodcentered{} CT 710 \textperiodcentered{} [7]}
A Bayesian Belief Network (BBN) is used in a medical diagnosis system. The network includes
the binary variables C (Patient has Cancer), T (Test result is Positive) and S (Patient is a
Smoker), with structure S $\rightarrow$ C $\rightarrow$ T and conditional probabilities
P(S=Yes)=0.3, P(C=Yes\,$|$\,S=Yes)=0.2, P(C=Yes\,$|$\,S=No)=0.05, P(T=Yes\,$|$\,C=Yes)=0.9,
P(T=Yes\,$|$\,C=No)=0.1. Explain how the Bayesian Belief Network enables reasoning under
uncertainty in this medical context.
\end{asked}

\sbsr{0.52}{%
\lead{The three kinds of reasoning, each with its number}
\begin{itemize}
  \item \textbf{Causal, top-down} $-$ \emph{a smoker walks in; how likely is a positive
        test?}
        \[ P(T|S) = 0.9(0.2) + 0.1(0.8) = \mathbf{0.26} \]
        against $P(T|\neg S) = 0.9(0.05) + 0.1(0.95) = 0.14$.
  \item \textbf{Prior belief}, before anyone is seen:
        \[ P(C) = 0.2(0.3) + 0.05(0.7) = 0.06 + 0.035 = \mathbf{0.095} \]
        \[ P(T) = 0.9(0.095) + 0.1(0.905) = 0.0855 + 0.0905 = \mathbf{0.176} \]
  \item \textbf{Diagnostic, bottom-up} $-$ \emph{the test came back positive; how likely
        is cancer?} This is Bayes on the second edge:
        \[ P(C|T) = \frac{0.9 \times 0.095}{0.176} = \frac{0.0855}{0.176}
           = \mathbf{0.4858} \]
  \item \textbf{Through the chain} $-$ \emph{the test is positive; how likely is the
        patient a smoker?}
        \[ P(S|T) = \frac{P(T|S)P(S)}{P(T)} = \frac{0.26 \times 0.3}{0.176}
           = \frac{0.078}{0.176} = \mathbf{0.4432} \]
\end{itemize}
{\color{sub}\footnotesize Consistency check: $P(T)$ computed through $C$ gives 0.176, and
computed through $S$ gives $0.26(0.3) + 0.14(0.7) = 0.078 + 0.098 = 0.176$. \checkmark \;
\mk{Two routes, one answer} $-$ that is the network being coherent, and
it is worth showing.}}{%
\lead{\mk{What to say in words}, which is what the question asks}
\begin{itemize}
  \item The network \textbf{separates what is known from what is asked}. The doctor can
        only measure causal probabilities $-$ how often smokers get cancer, how often
        cancer shows on a test $-$ and the network stores exactly those. Bayes then runs
        them backwards on demand.
  \item It handles \textbf{partial evidence}: with no observation the belief in cancer is
        9.5\,\%; knowing the patient smokes raises it to 20\,\%; a positive test raises it
        to \mk{48.6\,\%}. Each piece of evidence revises, none replaces.
  \item \textbf{Conditional independence keeps it small.} $T$ depends on $C$ alone: once
        the cancer status is known, smoking tells you nothing more about the test. So
        three variables need \textbf{5 numbers} instead of $2^3 - 1 = 7$, and the saving
        grows exponentially with the network.
  \item \textbf{It never claims certainty.} A positive test leaves the patient more likely
        \emph{not} to have cancer than to have it $-$ 51.4\,\% against 48.6\,\% $-$ which
        is precisely the kind of conclusion a logic-based KB cannot express at all.
\end{itemize}}

\penalty0\vspace{2pt}

\creamq{B12. The three papers that ask for an example of your own}{\m{7}\m{5+3}\m{3+5}}
{\color{sub}\footnotesize \yr{\textbf{73 Bh} \m{7}, \texttt{80 Ba} \m{5+3},
\textbf{80 Ch} \m{3+5}} \gap$\cdot$\gap no data printed, so
\mk{bring your own} $-$ and the smallest one that shows everything is the
two-node version of B9}

\sbs{%
\lead{The example to memorise}
\begin{itemize}
  \item \textbf{Network}: $Cancer \rightarrow Test$, two nodes, one edge.
  \item \textbf{CPTs}: $P(C) = 0.01$; \; $P(+|C) = 0.9$, \; $P(+|\neg C) = 0.08$.
  \item \textbf{Causal use} (down the arrow): a patient with cancer tests positive
        90\,\% of the time. Straight from the table, no calculation.
  \item \textbf{Diagnostic use} (against the arrow): a positive test means cancer with
        probability
        \[ \frac{0.9 \times 0.01}{0.9(0.01) + 0.08(0.99)} = \frac{0.009}{0.0882}
           = 10.2\,\% \]
  \item \mk{Three numbers, four lines, and it demonstrates the entire
        connection} the three papers ask about.
\end{itemize}}{%
\lead{The connection, stated}
\begin{itemize}
  \item \textbf{The network supplies the numbers Bayes needs.} $P(C)$ is the prior sitting
        at the root; $P(+|C)$ and $P(+|\neg C)$ are the CPT of the child; the denominator
        is the law of total probability over that same CPT.
  \item \textbf{Bayes supplies the direction the network lacks.} The arrows point from
        cause to effect, which is the direction knowledge is available in. Diagnosis needs
        the other direction, and Bayes is the only thing that turns it round.
  \item So the answer to \emph{``is there any connection''} is:
        \mk{a belief network is a factored store of exactly the
        conditional probabilities Bayes' theorem consumes}, and inference in the
        network \emph{is} repeated application of Bayes plus the product rule.
  \item \yr{\texttt{80 Ba}}'s ``how can inference be done'' is answered by the four modes
        in \S4.7: causal, diagnostic, intercausal and mixed. Name all four; work one.
\end{itemize}}

%=======================================================================
\T{How every answer above was checked}

\sbs{%
\lead{The proofs}
\begin{itemize}
  \item \code{verify.py} carries a small \textbf{first-order resolution prover}
        (unification with occurs check, factoring, set of support) and runs it on
        \mk{every clause set printed above}, in the exact form printed.
  \item It confirms the empty clause is derivable for all of them $-$ and, just as
        importantly, that it is \textbf{not} derivable for the three papers whose printed
        goal does not follow:
  \begin{itemize}
    \item \yr{\textit{72 Ma}} \emph{Shailendri \textbf{likes} chicken} $-$ not derivable;
          \emph{eats} is.
    \item \yr{\textbf{\texttt{82 Bh}}} \emph{Mr.\,A is a criminal} $-$ not derivable;
          $Criminal(Mr.R)$ is.
    \item \yr{79 Ash} \emph{Steve likes BK301} $-$ not derivable from
          \emph{likes $\rightarrow$ easy}; derivable from the converse.
  \end{itemize}
  \item It also confirms that \yr{\textbf{69 Bh}} and \yr{75 Ba} \textbf{cannot} be closed
        without the two bridging axioms this document adds, and that they close once those
        are supplied.
\end{itemize}}{%
\lead{The numbers}
\begin{itemize}
  \item Every Bayes posterior, every total-probability denominator and both network
        products are recomputed from the paper's own figures and asserted against the
        boxed value above. \mk{Nothing here was copied from a source.}
  \item Three of those checks are consistency tests rather than answers: the measles
        posteriors summing to 1, the three factory posteriors summing to 1, and
        \yr{\textbf{\texttt{82 Bh}}}'s $P(T)$ computed by two different routes through the
        network agreeing at 0.176.
  \item \mk{Run it before trusting a rebuild}:
\end{itemize}
\vspace{1pt}
\code{cd AI/ExamNotes/src \&\& python verify.py}
\vspace{2pt}
\begin{itemize}
  \item A textbook error that survives a check is still an error. The three
        \emph{Insights} corrections in this chapter $-$ the split of a negated
        conjunction, the free variable where a Skolem function belongs, and the five-step
        CNF list $-$ were each found by hand and then confirmed here.
\end{itemize}}

